% concatenated from main.tex
% SIAM Article Template
\documentclass[hidelinks,onefignum,onetabnum]{siamart250211}

% Information that is shared between the article and the supplement
% (title and author information, macros, packages, etc.) goes into
% ex_shared.tex. If there is no supplement, this file can be included
% directly.

% Packages and macros go here
\usepackage{lipsum}
\usepackage{amsfonts}
\usepackage{graphicx}
\usepackage{tabularx}
\usepackage{epstopdf}
\usepackage{algorithmic}
\usepackage{xcolor}

\usepackage{subcaption}     % Required for creating subfigures

\ifpdf
  \DeclareGraphicsExtensions{.eps,.pdf,.png,.jpg}
\else
  \DeclareGraphicsExtensions{.eps}
\fi

% Add a serial/Oxford comma by default.
\newcommand{\creflastconjunction}{, and~}

% Used for creating new theorem and remark environments
\newsiamremark{remark}{Remark}
\newsiamremark{hypothesis}{Hypothesis}
\crefname{hypothesis}{Hypothesis}{Hypotheses}
\newsiamthm{claim}{Claim}
\newsiamthm{assumption}{Assumption}
\newsiamremark{fact}{Fact}
\newsiamremark{example}{Example}
\crefname{fact}{Fact}{Facts}

\newcommand{\red}[1]{{\color{red} #1}}
\newcommand{\ii}{\mathrm{i}}
\newcommand{\ca}{\mathcal{A}}
\newcommand{\cc}{\mathcal{C}} % replaces \cc from the letter document style
\newcommand{\ce}{\mathcal{E}}
\newcommand{\ck}{\mathcal{K}}
\newcommand{\cm}{\mathcal{M}}
\newcommand{\cz}{\mathcal{Z}}
\newcommand{\tmod}{\ \mathsf{mod}\ }

\newcommand{\real}{\mathbb{R}}
\newcommand{\ints}{\mathbb{Z}}
\newcommand{\natu}{\mathbb{N}}
\newcommand{\ftwo}{\mathbb{F}_2}
\newcommand{\caln}{\mathcal{N}}

\newcommand{\err}{\mathrm{Err}}

\newcommand{\bsa}{\boldsymbol{a}}
\newcommand{\bsb}{\boldsymbol{b}}
\newcommand{\bsc}{\boldsymbol{c}}
\newcommand{\bsd}{\boldsymbol{d}}
\newcommand{\bse}{\boldsymbol{e}}
\newcommand{\bsg}{\boldsymbol{g}}
\newcommand{\bsk}{\boldsymbol{k}}
\newcommand{\bscn}{\boldsymbol{N}}
\newcommand{\bsq}{\boldsymbol{q}}
\newcommand{\bsu}{\boldsymbol{u}}
\newcommand{\bsv}{\boldsymbol{v}}
\newcommand{\bsw}{\boldsymbol{w}}
\newcommand{\bsx}{\boldsymbol{x}}
\newcommand{\bsy}{\boldsymbol{y}}
\newcommand{\bst}{\boldsymbol{t}}
\newcommand{\bsi}{\boldsymbol{i}}
\newcommand{\bsz}{\boldsymbol{z}}
\newcommand{\bszero}{\boldsymbol{0}}
\newcommand{\bsone}{\boldsymbol{1}}
\newcommand{\bsell}{{\boldsymbol{\ell}}}
\newcommand{\bsalpha}{\boldsymbol{\alpha}}
\newcommand{\bscalE}{\boldsymbol{\mathcal{E}}}
\newcommand{\bsM}{\boldsymbol{M}}
\newcommand{\bskappa}{\boldsymbol{\kappa}}
\newcommand{\bsgamma}{\boldsymbol{\gamma}}

\newcommand{\rd}{\,\mathrm{d}}
\newcommand{\mrd}{\mathrm{d}}
\newcommand{\dunif}{\mathbb{U}}
\newcommand{\mc}{\mathrm{MC}}
\newcommand{\qmc}{\mathrm{QMC}}
\newcommand{\rqmc}{\mathrm{RQMC}}
\newcommand{\simiid}{\stackrel{\mathrm{iid}}\sim}
\newcommand{\tod}{\stackrel{\mathrm{d}}\to}
\newcommand{\eqd}{\stackrel{\mathrm{d}}=}


\newcommand{\e}{\mathbb{E}}
\newcommand{\var}{\mathrm{Var}}
\newcommand{\cov}{\mathrm{Cov}}
\newcommand{\vol}{\mathrm{vol}}
\newcommand{\rank}{\mathrm{rank}}
\newcommand{\row}{\mathrm{Row}}

\newcommand{\one}{\mathbf{1}}
\newcommand{\bvhk}{\mathrm{BVHK}}
\newcommand{\dnorm}{\mathcal{N}}
\newcommand{\norm}[1]{\Vert #1 \Vert_1}

\newcommand{\giv}{\!\mid\!}
\newcommand{\phm}{\phantom{-}}
\newcommand{\tran}{\mathsf{T}}
\newcommand{\phe}{\phantom{=}}
\newcommand{\lset}{\mathcal{L}}

\newcommand{\walk}{\mathrm{wal}_k}
\newcommand{\walkj}{\mathrm{wal}_{k_j}}
\newcommand{\walbsk}{\mathrm{wal}_{\bsk}}
\newcommand{\bwalk}{_b\mathrm{wal}_k}
\newcommand{\walkappa}[1]{{\mathrm{wal}_{#1}}}

\newcommand{\med}{\mathrm{med}}
\newcommand{\daw}{\mathrm{Daw}}
\newcommand{\ch}{\mathcal{H}}
\newcommand{\calr}{\mathcal{R}}
\newcommand{\cp}{\mathcal{P}}
\newcommand{\hk}{\mathrm{HK}}
\newcommand{\re}{\mathrm{Re}}

\newcommand{\supp}{\boldsymbol{s}}
\newcommand{\tc}{\Tilde{c}}
\newcommand{\teta}{\eta}
\newcommand{\mix}{\mathrm{mix}}
\newcommand{\itr}{\mathbb{I}}
\newcommand{\icomp}{\mathtt{i}}

%define your comment here
\def\Du#1{{\textcolor{red}{Du: #1}}}
\def\Pan#1{{\textcolor{blue}{Pan: #1}}}
% Sets running headers as well as PDF title and authors
\headers{QMC importance sampling}{Z. Pan, D. Ouyang, and Z. He}

% Title. If the supplement option is on, then "Supplementary Material"
% is automatically inserted before the title.
\title{Quasi-Monte Carlo integration over $\mathbb{R}^s$ with boundary-damping importance sampling \thanks{Submitted to the editors DATE.
\funding{This work of the first author was funded by the Austrian Science Fund (FWF) Project DOI 10.55776/P34808. The work of the third author was funded by the Guangdong Basic and Applied Basic Research Foundation grant  2024A1515011876 and 2025A1515011888.}}}

% Authors: full names plus addresses.
\author{Zexin Pan\thanks{Johann Radon Institute for Computational and Applied Mathematics, Linz, Austria 
  (\email{zexin.pan@oeaw.ac.at}).} \and Du Ouyang\thanks{Department of Mathematical Sciences, Tsinghua University, Beijing 100084, People's Republic of China (\email{oyd21@mails.tsinghua.edu.cn}).} \and Zhijian He\thanks{Corresponding author. School of Mathematics, South China University of Technology, Guangzhou 510641, People's Repulic of China (\email{hezhijian@scut.edu.cn}).}
  }

\usepackage{amsopn}
\DeclareMathOperator{\diag}{diag}



% FundRef data to be entered by SIAM
%<funding-group specific-use="FundRef">
%<award-group>
%<funding-source>
%<named-content content-type="funder-name"> 
%</named-content> 
%<named-content content-type="funder-identifier"> 
%</named-content>
%</funding-source>
%<award-id> </award-id>
%</award-group>
%</funding-group>

\begin{document}

\maketitle

% REQUIRED
\begin{abstract}
This paper proposes a new importance sampling (IS) that is tailored to quasi-Monte Carlo (QMC) integration over $\mathbb{R}^s$. IS introduces a multiplicative adjustment to the integrand by compensating the sampling from the proposal instead of the target distribution. Improper proposals result in severe adjustment factor for QMC. Our strategy is to first design a adjustment factor to meet desired regularities and then determine a tractable transport map from the standard uniforms to the proposal for using QMC quadrature points as inputs.
The transport map has the effect of damping the boundary growth of the resulting integrand so that the effectiveness of QMC can be reclaimed. Under certain conditions on the original integrand, our proposed IS  enjoys a fast convergence rate independently of the dimension $s$, making it amenable to high-dimensional problems.
\end{abstract}

% REQUIRED
\begin{keywords}
Quasi-Monte Carlo, importance sampling, transport maps
\end{keywords}

% REQUIRED
\begin{MSCcodes}
41A63, 65D30, 97N40
\end{MSCcodes}

\section{Introduction}
Quasi-Monte Carlo (QMC) has gained its success in many fields, including computational finance \cite{lecuyer:2009,zhang:2021}, uncertainty quantification \cite{kaarnioja:2022,schillings:2020}. Although it has the potential to improve the convergence rate of plain Monte Carlo, the gain of QMC depends on the regularity of the underlying functions and the sampling proposals. Importance sampling (IS) provides a way to choose tractable sampling proposals instead of the underlying distribution, which is a widely used variance reduction technique in the Monte Carlo setting \cite{owen:2013}. IS is more than just a variance reduction method. Particularly, it is used within Bayesian statistics and Bayesian inverse problems as an approximation of the target distribution by weighted samples that are generated from some typical proposals \cite{agapiou:2017}. Recently, IS is combined with QMC to achieve a faster rate of convergence \cite{he:2023,ouyang:2024,zhang:2021}. IS can bring enormous gains in QMC, and it can also backfire, yielding a severe integrand with superfast growth boundary \cite{ouyang:2024} when simple QMC would have had a more regularity. It is an open problem to choose proposals in IS to  optimize the performance of QMC. In this paper, we propose a new IS to dampen the growth of the integrand around the boundary of the domain, which is preferable for QMC
integration of unbound integrands \cite{he:2023,haltavoid}. 

Throughout this paper, we let $\real=(-\infty,\infty)$ and $\itr=(0,1)$. We use $\bsx$ for coordinates in $\real^s$ and $\bsu$ for coordinates in $\itr^s$. Consider integral of the form
$$\mu=\int_{\real^s} f(\bsx) \prod_{j=1}^s \varphi(x_j)\rd \bsx,$$
where $f:\real^s \to \real$ is a real-valued integrand and $\varphi:\real\to \real$ is a probability density function. Let $\Phi(x)=\int_{-\infty}^x \varphi (y) \rd y$ be the cumulative distribution function (CDF) and $\Phi^{-1}:\itr \to \real$ be the inverse CDF (quantile function). QMC integration takes quadrature points in the unit cube. Before using QMC, the sampling proposal needs to be expressed as a transformation of the standard uniform distribution, say $T:\itr^d\to \real^s$. The transformation $T$ is also called a transport map in the literature \cite{liu2024transportquasimontecarlo}. The dimension $d$ does not necessarily agree with $s$, but for simplicity we take $d=s$ in this paper.  To estimate $\mu$, we consider estimators of the form
\begin{equation}\label{eqn:muhat}
 \hat{\mu}=\frac{1}{n}\sum_{i=0}^{n-1} w(\bsu_i)  f\circ T(\bsu_i) \quad \text{for} \quad T(\bsu)=(T_1(u_1),\dots, T_s(u_s)),   
\end{equation}
where $\{\bsu_0,\dots,\bsu_{n-1}\} \subseteq \itr^s$ is a deterministic QMC point set or randomized QMC (RQMC) point set for easy of error estimation \cite{vLEC02a,owen:2013}, the transport map $T_j:\itr\to \real$ is differentiable and the \textit{weight function}
$$w(\bsu)=\prod_{j=1}^s w_j(u_j)  \quad \text{for} \quad w_j(u)=T_j'(u)  \varphi\circ T_j(u) .$$
IS also takes the form \eqref{eqn:muhat}, in which the sampling proposal has independent marginals $T(u_j)$ and $w(\bsu)$ is  known as likelihood ratio. Many QMC methods proposed in the literature share this form. Examples are:
\begin{itemize}
    \item $T_j(u)=\Phi^{-1}(u)$, commonly known as inversion methods \cite{dick2025quasi, nichols:2014, haltavoid, ye2025medianqmcmethodunbounded}. For this case, $w(\bsu)=1$.
    \item $T_j(u)=a_j u+ b_j(1-u)$ for $a_j,b_j\in \real$, commonly known as truncation methods \cite{dick2018optimal, goda2024randomizing, kazashi2023suboptimality, nuyens2023scaled}.
    \item $T_j(u)=\Phi^{-1}_\nu(u)$ with $\Phi^{-1}_\nu(u)$ the inverse CDF of a Student's $t$-distribution with $\nu$ degree of freedom \cite{ouyang:2024}. A recent paper \cite{suzuki2025mobius} considers the M{\"o}bius transformation $T_j(u)=-\cot(\pi u)$, which can be viewed as a special case since $\cot(\pi u)$ is the inverse CDF of a Cauchy distribution.
\end{itemize}
We also note that there are many interesting methods beyond the above form \cite{dick2011quasi, dung2024optimal, he:2023, liu2024transportquasimontecarlo, wang2025convergence}.

While previous studies focus on the choice of $T_j$ and the resulting smoothness of $w_j$, we take a different perspective: we first design $w_j$ with required smoothness and then derive $T_j$ by solving the differential equation $w_j(u)=T_j'(u) \varphi\circ T_j(u)$, yielding
\begin{equation}\label{eqn:Tj}
    T_j(u)=\Phi^{-1}\Big(\int_{0}^u w_j(t)\rd t\Big).
\end{equation}
Specifically, we choose $w_j$ from a one-parameter family as follows. Let $\eta:[0,1/2]\to [0,1/2]$ be a differentiable monotonic function satisfying $\eta(0)=0$ and $\eta(1/2)=1/2$. For $\theta\in (0, 1/2]$, we define
\begin{equation}\label{eqn:wtheta}
w_\theta(u)=\begin{cases}
    (1-\theta)^{-1}\eta(u/\theta), &\text{ if } u\in (0,\theta/2] \\
    (1-\theta)^{-1}(1-\eta(1-u/\theta)) , &\text{ if } u\in (\theta/2,\theta) \\
    (1-\theta)^{-1}, &\text{ if } u\in [\theta,1/2] \\
    w_\theta(1-u), &\text{ if } u\in (1/2,1), \\
\end{cases}  
\end{equation}
where 
\begin{equation}\label{eq:tildeetap}
\eta(u)=\eta_p(u)=\begin{cases}
    2^{-p-2} u^{-p-1}\exp(2^p-u^{-p}), &\text{ if } u\in (0,1/2]\\
    0, &\text{ if } u=0
\end{cases} \text{ for } \ p\geq 1.
\end{equation}
While there are many admissible choices for $\eta$, we choose the function \eqref{eq:tildeetap} for its simplicity and computational convenience. By construction, $w_\theta(u)=w_\theta(1-u)$, $\sup_{u\in \itr} w_\theta(u)= (1-\theta)^{-1}$ and $\int_{0}^1 w_\theta(u)\rd u=1$.
Figure \ref{fig:wtheta} shows the case of $w_\theta(u)$ with $\theta = 0.2$ and $p=1$, which decays to zero quickly as $u$ approaches $0$ or $1$.  

\begin{figure}[ht]
    \centering
    \includegraphics[width=0.6\linewidth]{figure/wtheta.pdf}
    \caption{An illustration of $w_\theta(u)$ with $\theta =0.2$ and $p=1$.}
    \label{fig:wtheta}
\end{figure}

Given $\{\theta_j\mid j\in 1{:}s\}\subseteq (0,1/2]$ and $p\ge 1$, we choose $w_j(u)=w_{\theta_j}(u)$ for all $j=1,\dots,s$. Because $T_j$ maps $\itr$ onto $\real$, 
\begin{equation}\label{eqn:intfw}
    \int_{\itr^s}w(\bsu)f\circ T(\bsu) \rd\bsu=\int_{\real^s} f(\bsx) \prod_{j=1}^s \varphi(x_j)\rd \bsx=\mu.
\end{equation}
Hence, estimating $\mu$ is equivalent to integrating a new integrand $f^w(\bsu):=w(\bsu) f\circ T(\bsu)$
over $\itr^s$. The weight function $w(\bsu)$ can  dampen the growth of the integrand $f^w(\bsu)$ around the boundary of $\itr^s$. The proposed method is called the boundary-damping IS. Under appropriate choice of $\theta_j$, we show that $f^w$ has sufficient smoothness and can be efficiently integrated even in high-dimensional settings. Particularly, the boundary-damping IS yields a dimension-independent mean squared error rate under certain conditions on the integrand.


Our main contribution is three-fold. Firstly, we provide a novel IS to reclaim the performance of QMC by damping the boundary growth of the integrand. Unlike Laplace-based IS \cite{schillings:2020} that uses Gaussian proposals, our proposal is not a commonly used distribution but is easy to simulate. Secondly, we provide a rigorous error analysis based on generalized Fourier coefficients with an application to scrambled digital nets \cite{rtms}.   Thirdly, under certain conditions on the function $f(\bsx)$, our proposed IS can enjoy a faster convergence rate  than Monte Carlo while breaking the curse of dimensionality, making it amenable to high-dimensional problems. Lattice-based QMC quadrature rules can be constructed to yield asymptotic error bound independently of the dimension for unbounded integrands in weighted reproducing kernel Hilbert spaces with POD (``product and order dependent") weights \cite{nichols:2014}. Our analysis does not need to introduce weighted spaces and the digital net quadrature points are off-the-shelf.

The rest of this paper is organized as follow. Section~\ref{sec:back} provides some preliminaries on function norms and digital nets. Section~\ref{sec:bounds} studies upper bounds on generalized Fourier coefficients of the integrand. Section~\ref{sec:norm} investigates the norms of $f$ and $f^w$. Section~\ref{sec:main} conducts the numerical analysis on scrambled net-based integration. Numerical results are presented in Section~\ref{sec:numer} to show the effectiveness of our proposed method. Section~\ref{sec:conclusion} concludes this paper.



\section{Preliminaries}\label{sec:back}
We first introduce some notations that will be used in the following. For a vector $\bsx\in\real^s$ and a subset $v\subseteq 1{:}s$, $\bsx_v$  denotes the subvector of $\bsx$ indexed by $v$, while $\bsx_{v^c}$ denotes the subvector indexed by $1{:}s\setminus v$. Let $\natu$ be the set of positive integers and $\natu_0 = \natu\cup\{0\}$. For $\bsk\in \natu_0^s$, $\supp(\bsk):=\{j\in 1{:}s\mid k_j\neq 0\}$ denotes the support of $\bsk$. Let $\Vert\bsx\Vert_q=(\sum_{j=1}^s |x_j|^q)^{1/q}$ for $q>0$. For a nonempty set $v\subseteq 1{:}s$, define the mixed derivative by
$$\partial^v f(\bsx) = \left(\prod_{j\in v}\frac{\partial}{\partial x_j}\right)f(\bsx),$$
and define $\partial^\emptyset f(\bsx)=f(\bsx)$ by convention. Let $\mathbb{C}$ be the set of complex numbers. 
We define the operator
$\mathcal{I}:L^\infty(\itr)\to L^\infty(\itr)$ by
$$\mathcal{I}(\phi)(u)=\int_0^u \phi(t) \rd t.$$ 
All constants in this paper have an implicit dependency on $\varphi$ and we suppress it from the notation for simplicity.


\subsection{Function norms}
We assume that the density $\varphi (x)$  is a strictly  positive, bounded, symmetric, light-tailed  function in the following.

\begin{assumption}~\label{assump:rho}
Assume that $\varphi(x)>0$, $\varphi(x)=\varphi(-x)$, $\varphi_\infty:=\sup_{x\in \real} \varphi(x)<\infty$, and for any $\varepsilon>0$, there exists a constant $c_\varepsilon>0$ such that
$\varphi(x)\geq c_\varepsilon \Phi(x)^{1+\varepsilon}$ for any $x\le 0$. 
\end{assumption}

Under Assumption~\ref{assump:rho},  the CDF $\Phi(x)$ and  its inverse $\Phi^{-1}(u)$ is differentiable and strictly increasing. It is easy to verify the Gaussian density $\varphi(x)=\exp(-x^2/2)/\sqrt{2\pi}$ satisfies Assumption~\ref{assump:rho} due to the inequality \cite{gordon1941values} 
$$\int_x^\infty \varphi(y)\rd y \leq \frac{1}{x}\varphi(x) \text{ for } x> 0.$$
The condition on symmetry can be dropped by imposing a proper condition on the right tail of the density $\varphi(x)$. For simplicity, we work on symmetrical densities.

This paper will use some function norms as defined in the following:
\begin{align*}
    \Vert f \Vert_{L^q(\itr^s)}&=\begin{cases}
        \Big(\int_{\itr^s} |f(\bsu)|^q \rd \bsu \Big)^{1/q}&0<q<\infty\\
        \sup_{u\in \itr^s}|f(\bsu)|&q=\infty
    \end{cases},\\
    \Vert f \Vert_{L^q(\real^s,\varphi)}&=\Big(\int_{\real^s} |f(\bsx)|^q  \prod_{j=1}^s \varphi(x_j) \rd\bsx\Big)^{1/q},\\
    \Vert f \Vert_{W^{1,q}_{\mix}(\real^s,\varphi)}&=\Big(\sum_{v\subseteq 1{:}s} \Vert \partial^v  f \Vert^q_{L^q(\real^s,\varphi)}\Big)^{1/q},\\
    \Vert f \Vert_{L^{q,\infty}(\real^s,\varphi)}&=\Big(\sup_{\bsx\in \real^s} |f(\bsx)|^q  \prod_{j=1}^s \frac{\varphi(x_j)}{\varphi_\infty}  \Big)^{1/q},\\
    \Vert f \Vert_{W^{1,q,\infty}_{\mix}(\real^s,\varphi)}&=\Big(\sum_{v\subseteq 1{:}s} \Vert \partial^v  f \Vert^q_{L^{q,\infty}(\real^s,\varphi)}\Big)^{1/q}.
\end{align*}
Notice that for constant functions $f(\bsx)=c$, all of the above norms reduce to $|c|$. More generally, if $f(\bsx)$ does not depend on $x_{v^c}$ for $v\subseteq 1{:}s$, we can treat $f$ either as a function over $\real^s$ or a function over $\real^{|v|}$ by removing redundant input variables $x_{v^c}$. Under either interpretation, the resulting norm is the same according to the definitions given above. This property will be useful later when we analyze the ANOVA decomposition of $f$. The next lemma shows $f\in L^{q,\infty}(\real^s,\varphi)$ almost implies $f\in L^{q}(\real^s,\varphi)$.



\begin{lemma}\label{lem:embedding}
    For $\varphi$ satisfying Assumption~\ref{assump:rho} and $q>q'\geq 1$,
    $$\Vert f \Vert_{L^{q'}(\real^s,\varphi)}\leq C^s_{q,q'}\Vert f\Vert_{L^{q,\infty}(\real^s,\varphi)}$$
    and 
    $$\Vert f \Vert_{W^{1,q'}_{\mix}(\real^s,\varphi)}\leq \tilde C^s_{q,q'}\Vert f\Vert_{W^{1,q,\infty}_{\mix}(\real^s,\varphi)},$$
    where  $C_{q,q'}, \tilde C_{q,q'}$ are constants depending on $q$ and $q'$.
\end{lemma}
\begin{proof}
    See the appendix.
\end{proof}

\begin{remark}\label{rem:rqmcrate}
   If $f\in W^{1,q,\infty}_{\mix}(\real^s,\varphi)$ and $\varphi(x)=\exp(-x^2/2)/\sqrt{2\pi}$, we have for any $v\subseteq1{:}s$,
   $$|\partial^v  f(\bsx)|\le c \left(\prod_{j=1}^s \varphi(x_j)\right)^{-1/q}=c'\exp\left(\frac{1}{2q}\Vert\bsx\Vert_2^{2}\right)$$
   for some $c,c'>0$. This implies that $f(\bsx)$ belongs to the superfast growth class $G_e(1/(2q),c',2)$ defined in \cite{ouyang:2024}. For $q>1$, RQMC integration with the usual inversion method yields a root mean squared error of $O(n^{-1+1/q+\epsilon})$ for arbitrarily small $\epsilon>0$ \cite{he:2023,haltavoid}, where the implied constant depends on the dimension $s$ and $\epsilon$.
\end{remark}


\subsection{ANOVA decomposition over $\real^s$}

 Following the framework in \cite{kuo2010decompositions}, we introduce the generalized ANOVA decomposition. Let $P_j:L^1(\real^s,\varphi)\to L^1(\real^s,\varphi)$ denote the integration operator
$$P_j(f)(\bsx)=\int_{\real} f(\bsx) \varphi(x_j) \rd x_j \text{ for } \bsx \in \real^{s}.$$
Notice that $P_j(f)$ does not depend on $x_j$ and $P^2_j=P_j$. We further define the iterated integration operator $P_v=\prod_{j\in v}P_j$. By Fubini's theorem, $P_j$ in $P_v$ can be applied in any order. By convention, $P_\emptyset=I$ is the identity operator. The ANOVA decomposition of $f\in L^1(\real^s,\varphi)$  is given by
\begin{equation}\label{eqn:anovadecomp}
 f=\sum_{v\subseteq 1{:}s} f_v \text{ for } f_v=\Big(\prod_{j\in v} (I-P_j)\Big) P_{1{:}s\setminus v} f.   
\end{equation}
It follows that if $j\in v$,
\begin{equation}\label{eqn:ANOVAPj}
  P_j (f_v)=(P_j-P_j^2)\Big(\prod_{j'\in v, j'\neq j} (I-P_{j'})\Big) P_{1{:}s\setminus v} f =0.  
\end{equation}
The next lemma shows each ANOVA component $f_v$ inherits the smoothness of $f$. 
\begin{lemma}\label{lem:anova}
    If $f\in W^{1,q}_{\mix}(\real^s,\varphi)$, then $f_v\in W^{1,q}_{\mix}(\real^s,\varphi)$ for all $v\subseteq 1{:}s$. Furthermore, if $\varphi$ satisfies Assumption~\ref{assump:rho} and $f\in W^{1,q,\infty}_{\mix}(\real^s,\varphi)$ with $q>1$, then $f_v\in W^{1,q,\infty}_{\mix}(\real^s,\varphi)$ for all $v\subseteq 1{:}s$.
\end{lemma}
\begin{proof}
    See the appendix.
\end{proof}

\subsection{Scrambled $(t,m,s)$-nets integration}\label{subsec:walsh} 
In this paper, we use scrambled $(t,m,s)$-nets in base $b\ge 2$ as the quadrature points in the estimator \eqref{eqn:muhat} as defined in the following.

\begin{definition}
For $t,\ m\in\natu_0$ and an integer $b\ge 2$, a  set $\cp:=\{\bsu_0,\dots,\bsu_{b^m-1}\}$ in $[0,1)^s$ is called
a $(t,m,s)$-net in base $b$ if every interval of the form $\prod_{j=1}^s\left[\frac{a_j}{b^{k_j}},\frac{a_j+1}{b^{k_j}}\right)$ contains exactly $b^t$ points of $\cp$ for all integers $a_j\in[0,b^{k_j})$ and all $k_j\in \natu_0$ satisfying $\sum_{j=1}^sk_j=m-t$. For $\emptyset \neq w\subseteq 1{:}s$, the projection of $\cp$ on  coordinates $j\in w$ forms a $(t_w,m,|w|)$-net in base $b$, where $t_\omega\in \natu$ is called $t$-quality parameter and $t_{w}=t$ when $w=1{:}s$. When performing Owen's scrambling \cite{rtms} on $\cp$, the resulting point set is called the scrambled $(t,m,s)$-net, which is also a $(t,m,s)$-net with probability one.
\end{definition}

\begin{remark}
As shown in \cite{c812a4a5-75ac-3614-936d-5782e44941d0}, for the Sobol' sequence \cite{sobol67} and the Niederreiter sequence \cite{NIEDERREITER198851}, the $t$-quality parameters $t_w$ satisfies that
\begin{equation}\label{eqn:tutj}
    t_\omega\leq \sum_{j\in \omega}t_j
\end{equation}
with $t_j= O(\log_b(j))$.
\end{remark}


%\subsection{Walsh functions and Walsh decomposition}\label{subsec:walsh}

Let $b$ be a prime number from now on. For $k\in \natu_0$ with $b$-adic expansion $k=\sum_{i=1}^r \kappa_i b^{i-1}$, the $k$-th $b$-adic Walsh function is given by
\begin{equation}\label{eqn:walk}
    {}_b\walk (u)=\exp\left(\frac{2\pi \ii}{b}\sum_{i=1}^r \kappa_i v_i\right),
\end{equation}
where $\ii=\sqrt{-1}$,  $v_1,\dots,v_{r}$ is given by the $b$-adic expansion $u=\sum_{i=1}^\infty v_i b^{-i}\in [0,1)$. For the multivariate case, the Walsh functions are defined by
$$
{}_b\walbsk (\bsu) :=\prod_{j=1}^s {}_b\walkj(u_j),\quad \bsk\in \natu_0^s.
$$
It is already known that $\{{}_b\walbsk (\bsu)\mid \bsk\in \natu_0^s\}$ forms an orthonormal basis of $L^2(\itr^s)$; see \cite{dick:pill:2010} for example. For $f\in L^2(\itr^s)$, we  have the Walsh series expanding
$$
f(\bsu)\sim \sum_{\bsk\in\natu^s} \hat{f}(\bsk){}_b\walbsk (\bsu),
$$
where $f\sim g$ denotes the equivalence relation for the $L^2(\itr^s)$ space, and $$\hat{f}(\bsk):=\int_{\itr^s}f(\bsu)\overline{{}_b\walbsk (\bsu)} \rd \bsu$$ denotes the $\bsk$-th Walsh coefficient of $f$. Let
$$\sigma^2_{\bsell}=\sum_{\bsk\in L_{\bsell}} |\hat{f}(\bsk)|^2,$$
where 
\begin{equation}\label{eqn:Ll}
    L_{\bsell}=\{\bsk\in \natu_0^s\mid \supp(\bsk)=\supp(\bsell), b^{\ell_j-1} \leq k_j < b^{\ell_j} \ \forall j\in \supp(\bsell)\}.
\end{equation}
The next lemma provides a useful upper bound for  the scrambled net variance.

\begin{lemma}\label{lem:qmcvar}
For $f\in L^2(\itr^s)$ and $\{\bsu_0,\dots,\bsu_{n-1}\}$ a scrambled  $(t,m,s)$-net in base $b\ge 2$ with $n=b^m$,
\begin{equation}\label{eqn:RQMCvar}
    \e\left|\frac{1}{n}\sum_{i=0}^{n-1} f(\bsu_i)-\int_{\itr^s} f(\bsu)\rd \bsu\right|^2=\frac{1}{n}\sum_{\emptyset\neq \omega \subseteq 1{:}s}\sum_{\bsell\in \natu^\omega} \Gamma_{\omega,\bsell} \sigma^2_{\bsell},
\end{equation}
    where $\natu^{v}=\{\bsell\in \natu^s_0\mid \supp(\bsell)=v\}$ and 
    \begin{equation}\label{eqn:gaincoefbound}
        \Gamma_{\omega,\bsell}\leq \Big(\frac{b}{b-1}\Big)^{|\omega|-1}b^{t_\omega}\bsone\{\Vert \bsell\Vert_1> m-t_\omega-|\omega|\}.
    \end{equation}
\end{lemma}
\begin{proof}
    Equation~\eqref{eqn:RQMCvar} is first derived by \cite{owen96,snxs} in terms of Haar wavelet basis. See \cite[Theorem 13.6]{dick:pill:2010} for the version using Walsh basis. Inequality~\eqref{eqn:gaincoefbound} comes from \cite{GODA2023101722}.
\end{proof}

\section{Generalized Fourier coefficients and their bounds}\label{sec:bounds}

Recall that our estimator $\hat{\mu}$ is given by Equation~\eqref{eqn:muhat}, where the integrand is $f^w(\bsu)=w(\bsu) f\circ T(\bsu)$.
A crucial step in analyzing the error of $\hat{\mu}$ is to bound the generalized Fourier coefficients of $f^w$. Specifically, we consider a complete orthonormal basis $\{\phi_k(u)\mid k\in \natu_0\}$ of $L^2(\itr)$ with $\phi_0(u)=1$. We additionally assume every $\phi_k \in L^\infty(\itr)$. Examples are trigonometric functions used in the analysis of lattice rules \cite{dick2022lattice} and Walsh functions introduced in Subsection~\ref{subsec:walsh}. Given $f\in L^1(\itr^s)$, we define the generalized Fourier coefficients
$$\hat{f}(\bsk)=\int_{\itr^s} f(\bsu) \prod_{j=1}^s \overline{\phi_{k_j}(u_j)}\rd \bsu =\int_{\itr^s} f(\bsu)\prod_{j\in \supp(\bsk)} \overline{\phi_{k_j}(u_j)}\rd \bsu$$
for $\bsk=(k_1,\dots,k_s)\in \natu^s_0$.
Notice that $\{\hat{f}(\bsk)\mid \bsk \in \natu^s_0\}$ is not square-summable if $f\notin L^2(\itr^s)$.

 The next lemma shows $\hat{f}^w(\bsk)$ only depends on the ANOVA components $f_v$ with $v\subseteq \supp(\bsk)$.

\begin{lemma}\label{lem:fwk}
    For $f\in L^q(\real^s,\varphi)$ and $w_j=w_{\theta_j}$ defined by Equation~\eqref{eqn:wtheta} with  $\{\theta_j\mid j\in 1{:}s\}\subseteq (0,1/2]$, 
    \begin{equation}\label{eqn:fwLqbound}
        \Vert f^w\Vert_{L^q(\itr^s)}\leq 2^{\frac{q-1}{q}s}\Vert f \Vert_{L^q(\real^s,\varphi)}
    \end{equation}
    and     
    $$\hat{f}^w(\bsk)=\sum_{v\subseteq\supp(\bsk)}\hat{f}^w_v(\bsk_v)\prod_{j\in \supp(\bsk)\setminus v} \hat{w}_j(k_j),$$
where
   $$f^w_v(\bsu_v)=f_v\circ T(\bsu) \prod_{j\in v} w_j(u_j)$$
    with $f_v\circ T(\bsu)$ interpreted as a function of $\bsu_v$ by removing redundant variables $\bsu_{v^c}$.
\end{lemma}
\begin{proof}
Because $T'_j(u)\geq 0$ and $\sup_{u\in\itr} w_j(u)= (1-\theta_j)^{-1}\leq 2$
\begin{align*}
    \Vert f^w\Vert^q_{L^q(\itr^s)}=&\int_{\itr^s} |f|^q\circ T(\bsu)\prod_{j=1}^s w_j(u_j)^q\rd\bsu\\
    \leq &\prod_{j=1}^s\Big(\sup_{u\in\itr} w_j(u)\Big)^{q-1}\int_{\itr^s}|f|^q\circ T(\bsu)\prod_{j=1}^s\varphi\circ T_j(u_j) T_j'(u_j)\rd\bsu\\
    \le &2^{(q-1)s} \int_{\real^s} |f(\bsx)|^q \prod_{j=1}^s \varphi(x_j)\rd \bsx=2^{(q-1)s} \Vert f \Vert^q_{L^q(\real^s,\varphi)}.
\end{align*}
Next, Equation~\eqref{eqn:ANOVAPj} implies for $j\in v$, 
$$\int_{\itr}  f_v \circ T(\bsu) w_j(u_j)\rd u_j=\int_{\itr} f_v \circ T(\bsu) \varphi\circ T_j(u_j) T_j'(u_j) \rd u_j=\int_\real f_v(\bsx) \varphi(x_j)\rd x_j=0.$$
Together with $\int_{0}^1 w_\theta(u)\rd u=1$, we have
\begin{align*}
    \hat{f}^w(\bsk)=& \sum_{v\subseteq 1{:}s} \int_{\itr^s} f_v\circ T(\bsu) \prod_{j=1}^s w_j(u_j) \prod_{j\in \supp(\bsk)} \overline{\phi_{k_j}(u_j)}\rd \bsu \\
    %=& \sum_{v\subseteq \supp(\bsk)} \int_{\itr^{|\supp(\bsk)|}} \Big(\sum_{v'\subseteq \supp(\bsk)^c}\int_{\itr^{s-|\supp(\bsk)|}} f_{v\cup v'}\circ T(\bsu) \\
   % &\prod_{j\in  \supp(\bsk)^c } w_j(u_j) \rd \bsu_{ \supp(\bsk)^c} \Big)\prod_{j\in \supp(\bsk)} w_j(u_j)\overline{\phi_{k_j}(u_j)}\rd \bsu_{\supp(\bsk)}\\
   % =& \sum_{v\subseteq \supp(\bsk)} \int_{\itr^{|\supp(\bsk)|}}f_{v}\circ T(\bsu)\prod_{j\in \supp(\bsk)} w_j(u_j)\overline{\phi_{k_j}(u_j)}\rd \bsu_{\supp(\bsk)}\\
    =& \sum_{v\subseteq\supp(\bsk)}\int_{\itr^{|v|}}  f_v\circ T(\bsu) \prod_{j\in v} w_j(u_j)\overline{\phi_{k_j}(u_j)}\rd \bsu_v\prod_{j\in \supp(\bsk)\setminus v} \int_\itr w_j(u_j)\overline{\phi_{k_j}(u_j)}\rd u_j\\
    =&\sum_{v\subseteq\supp(\bsk)}\hat{f}^w_v(\bsk_v)\prod_{j\in \supp(\bsk)\setminus v} \hat{w}_j(k_j).
\end{align*}
\end{proof}

Lemma~\ref{lem:fwk} suggests $|\hat{f}^w(\bsk)|$ is controlled by $|\hat{w}_j(k_j)|$ for $j\in \supp(\bsk)$ and $|\hat{f}^w_v(\bsk_v)|$ for $v\subseteq \supp(\bsk)$. The next lemma bounds the generalized Fourier coefficients of $w_j=w_{\theta_j}$.

\begin{lemma}\label{lem:wkbound}
For $k\in \natu$ and $w_\theta$ with $\theta\in (0,1/2]$,  
$$|\hat{w}_\theta(k)|\leq  4\min\big(\theta\Vert\phi_k\Vert_{L^\infty(\itr)}, \Vert \mathcal{I}(\phi_k)\Vert_{L^\infty(\itr)}\big).$$
\end{lemma}
\begin{proof}
    Because $\phi_k(u)$ is orthogonal to $\phi_0(u)=1$,
    $$\mathcal{I}(\phi_k)(1)=\int_\itr \phi_k(u)\rd u =0.$$
    Since $0\leq w_\theta(u)\leq (1-\theta)^{-1}$,
    \begin{align*}
        |\hat{w}_\theta(k)|=&\Big|\int_\itr w_\theta(u) \overline{\phi_k(u)}\rd u\Big| = \Big|\int_\itr w_\theta(u) \phi_k(u)\rd u\Big|\\ =& \Big|\int_\itr \big(w_\theta(u)-(1-\theta)^{-1}\big) \phi_k(u)\rd u\Big|\\
        \leq & \Big|\int_0^\theta \big((1-\theta)^{-1}-w_\theta(u)\big) \phi_k(u)\rd u\Big| +  \Big|\int_{1-\theta}^1  \big((1-\theta)^{-1}-w_\theta(u)\big) \phi_k(u)\rd u\Big|\\
        \leq & 2\theta(1-\theta)^{-1}   \Vert\phi_k\Vert_{L^\infty(\itr)}.
    \end{align*}
    Next, we use integration by parts and $\mathcal{I}(\phi_k)(0)=\mathcal{I}(\phi_k)(1)=0$ to obtain
    \begin{align*}
|\hat{w}_\theta(k)|=&\Big|\int_\itr w_\theta(u) \phi_k(u)\rd u\Big|
=\Big|\int_\itr w'_\theta(u) \mathcal{I}(\phi_k)(u)\rd u\Big|\\
         \leq & \Big|\int_0^{\theta} w'_\theta(u) \mathcal{I}(\phi_k)(u)\rd u\Big|+\Big|\int_{1-\theta}^{1} w'_\theta(u) \mathcal{I}(\phi_k)(u)\rd u\Big|\\
         \leq &  2\Vert \mathcal{I}(\phi_k)\Vert_{L^\infty(\itr)}\int_0^{\theta} w'_\theta(u) \rd u\\
         =& 2 (1-\theta)^{-1} \Vert \mathcal{I}(\phi_k)\Vert_{L^\infty(\itr)}.
  \end{align*}
    The conclusion follows once we bound $(1-\theta)^{-1}\leq 2$.
\end{proof}

Bounds on $|\hat{f}^w_v(\bsk_v)|^2$ generally depend on the smoothness of $f^w_v$ and properties of the orthonormal basis. In this work, we use the following bound based on $\Vert\partial^{v}f_v^w \Vert_{L^2(\itr^{|v|})}$. Whether this norm is finite and how it depends on $\{\theta_j\mid j\in v\}$ is the subject of the next section.

\begin{lemma}\label{lem:fkbound}
Assume there exist $\lambda:\natu\to \mathbb{C}\setminus \{0\}$, $\mathcal{N}:\natu\to\natu$ and  $\mathcal{M}:\natu\to\natu_0$ such that
\begin{equation}\label{eqn:intphik}
    \mathcal{I}(\phi_k)=\lambda(k)\phi^*_{\mathcal{N}(k)}\phi_{\mathcal{M}(k)},
\end{equation}
where $\{\phi^*_{\ell}(u)\in L^\infty(\itr)\mid \ell\in \natu\}$ satisfies $\Vert\phi^*_{\ell}\Vert_{L^\infty(\itr)}=1$ for every $\ell\in \natu$. Then for $f\in W^{1,2}_{\mix}(\itr^s)$ and $\bsell\in \natu^s$,
    $$\sum_{\bsk\in \mathcal{N}_s^{-1}(\bsell) }  |\hat{f}(\bsk)|^2\prod_{j=1}^s |\lambda(k_j)|^{-2}\leq  \Vert \partial^{1{:}s}f \Vert^2_{L^2(\itr^s)},$$
    where $\mathcal{N}_s^{-1}(\bsell)=\{\bsk\in \natu^s\mid \mathcal{N}(k_j)=\ell_j, j\in 1{:}s\}.$
\end{lemma}
\begin{proof}
    Since $\mathcal{I}(\phi_{k_1})(0)=\mathcal{I}(\phi_{k_1})(1)=0$, we can use integration by parts for weak derivatives \cite[Chapter 4]{nikol2012approximation} and derive
    $$\hat{f}(\bsk)=\int_{\itr^s} f(\bsu) \prod_{j=1}^s \overline{\phi_{k_j}(u_j)}\rd \bsu =\int_{\itr^s}  \partial^{1} f(\bsu) \overline{\mathcal{I}(\phi_{k_1})(u_1)} \prod_{j=2}^s \overline{\phi_{k_j}(u_j)}\rd \bsu.$$
    After repeating the above calculation for $u_2,\dots,u_s$, we get
    \begin{align*}
        \hat{f}(\bsk)=&\int_{\itr^s} \partial^{1{:}s}f(\bsu) \prod_{j=1}^s \overline{\mathcal{I}(\phi_{k_j})(u_j)}\rd \bsu  \\
        = &\int_{\itr^s} \partial^{1{:}s}f(\bsu) \prod_{j=1}^s \overline{\lambda(k_j)\phi^*_{\mathcal{N}(k_j)}(u_j)\phi_{\mathcal{M}(k_j)}(u_j)}\rd \bsu.
    \end{align*}
    Over $\bsk\in \mathcal{N}_s^{-1}(\bsell)$, $\mathcal{N}(k_j)=\ell_j$. Moreover, if $\mathcal{M}(k_j)=\mathcal{M}(k'_j)$ for $k_j,k'_j\in \mathcal{N}^{-1}(\ell_j)$,  
    $$\mathcal{I}(\phi_{k_j})=\lambda(k_j)\phi^*_{\ell_j}\phi_{\mathcal{M}(k_j)}=\frac{\lambda(k_j)}{\lambda(k'_j)}\lambda(k'_j)\phi^*_{\ell_j}\phi_{\mathcal{M}(k'_j)}=\frac{\lambda(k_j)}{\lambda(k'_j)}\mathcal{I}(\phi_{k'_j}),$$
    which is true only if $k_j=k'_j$. Therefore, the mapping $\bsk\to (\mathcal{M}(k_1),\dots,\mathcal{M}(k_s))$ is injective over $\mathcal{N}_s^{-1}(\bsell)$ and
    \begin{align*}
       &\sum_{\bsk\in \mathcal{N}_s^{-1}(\bsell) }  |\hat{f}(\bsk)|^2\prod_{j=1}^s |\lambda(k_j)|^{-2}\\
       \leq & \sum_{\bsk\in \mathcal{N}_s^{-1}(\bsell) }  \Big|\int_{\itr^s} \Big(\partial^{1{:}s}f(\bsu) \prod_{j=1}^s \overline{\phi^*_{\ell_j}(u_j)}\Big) \prod_{j=1}^s\overline{\phi_{\mathcal{M}(k_j)}(u_j)}\rd \bsu\Big|^2\\
        \leq & \int_{\itr^s} \Big|\partial^{1{:}s}f(\bsu) \prod_{j=1}^s \overline{\phi^*_{\ell_j}(u_j)}\Big|^2\rd \bsu
        \leq  \Vert \partial^{1{:}s}f \Vert^2_{L^2(\itr^s)},
    \end{align*}
    where in the second inequality we have applied Bessel's inequality for the orthonormal basis $\{\prod_{j=1}^s\phi_{k_j}(u_j)\mid \bsk\in \natu^s_0\}$.
\end{proof}

\begin{example}\label{example}
    To illustrate, we prove Equation~\eqref{eqn:intphik} holds for  $b$-adic  Walsh functions. Let 
    $\phi_k(u)={}_b\walk (u)$ defined by Equation~\eqref{eqn:walk}.  For $k\in \natu$ satisfying $b^{r-1}\leq k<b^r$, because $u_1,\dots,u_{r-1}$ in the $b$-adic expansion $u=\sum_{i=1}^\infty u_i b^{-i}$ is constant over the interval $\itr_a=[a b^{-r+1},(a+1)b^{-r+1})$ with an integer $a\in [0,b^{r-1}-1]$, 
    \begin{align*}
        \int_{\itr_a} \phi_k(u) \rd u=\exp\Big(\frac{2\pi \ii}{b}\sum_{i=1}^{r-1} \kappa_i u_i\Big)\int_{\itr_a} \exp\Big(\frac{2\pi \ii}{b} \kappa_r u_{r}\Big) \rd u=0,
    \end{align*}
    where the sum from $1$ to $r-1$ is set to $0$ if $r=1$. Hence for $u\in \itr_a$, 
    \begin{align*}
       \mathcal{I}(\phi_k)(u)=
       &\exp\Big(\frac{2\pi \ii}{b}\sum_{i=1}^{r-1} \kappa_i u_i\Big) \int_{a b^{-r+1}}^{u}  \exp\Big(\frac{2\pi \ii}{b} \kappa_r u'_{r}\Big) \rd u'\\
       =&\exp\Big(\frac{2\pi \ii}{b}\sum_{i=1}^{r-1} \kappa_i u_i\Big) \int_{0}^{u}  \exp\Big(\frac{2\pi \ii}{b} \kappa_r u'_{r}\Big) \rd u'\\
       =&\phi_{\mathcal{M}(k)}(u)\mathcal{I}(\phi_{\mathcal{N}(k)})(u), 
    \end{align*}
    where $\mathcal{M}(k)=k-\kappa_r b^{r-1}$ and $\mathcal{N}(k)=\kappa_r b^{r-1}$. Equation~\eqref{eqn:intphik} follows after we write $\mathcal{I}(\phi_{\mathcal{N}(k)})=\lambda(k) \phi^*_{\mathcal{N}(k)}$ with 
    $$\lambda(k)=\Vert\mathcal{I}(\phi_{\mathcal{N}(k)})\Vert_{L^\infty(\itr)}=\sup_{u\in\itr}\Big|\int_{0}^{u}  \exp\Big(\frac{2\pi \ii}{b} \kappa_r u'_{r}\Big) \rd u'\Big|\leq b^{-r+1}.$$
\end{example}

Lemmas~\ref{lem:fwk}-\ref{lem:fkbound} together give the following bound on $|\hat{f}^w(\bsk)|^2$.
\begin{theorem}\label{thm:fkbound}
    Suppose $\{\phi_k(u)\mid k\in \natu_0\}$ satisfies Equation~\eqref{eqn:intphik} and 
    $$\sup_{k\in \natu_0} \Vert \phi_k\Vert_{L^\infty(\itr)} =C_\phi<\infty.$$
    Then for $\alpha\in (0,1)$, $\bsell\in \natu_0^{s}$ and $f^w\in W^{1,2}_{\mix}(\itr^s)$ with $w_j=w_{\theta_j}$ for $\{\theta_j\mid j\in 1{:}s\}\subseteq (0,1/2]$,
    \begin{multline}\label{eqn:fwkbound}
        \sum_{\bsk\in \mathcal{E}(\bsell) } |\hat{f}^w(\bsk)|^2\prod_{j\in\supp(\bsell)} |\lambda(k_j)|^{-2\alpha}\\
        \leq 2^{|\supp(\bsell)|}\sum_{v\subseteq\supp(\bsell)}\Vert \partial^{v}f^w_v \Vert^{2\alpha}_{L^2(\itr^{|v|})}\Vert f^w_v \Vert^{2-2\alpha}_{L^2(\itr^{|v|})} \prod_{j\in \supp(\bsell)\setminus v} \mathcal{S}(\theta_j,\ell_j,\alpha),
    \end{multline}
    where 
    $$\mathcal{E}(\bsell)=\Big\{\bsk\in \natu_0^{s}\mid \supp(\bsk)=\supp(\bsell), \mathcal{N}(k_j)=\ell_j \ \forall j\in \supp(\bsell)\Big\},$$
    $$\mathcal{S}(\theta,\ell,\alpha)=16C^2_\phi\sum_{k\in\mathcal{N}^{-1}(\ell) } \min\big(\theta^2 , |\lambda(k)|^{2}\big)|\lambda(k)|^{-2\alpha}.$$
\end{theorem}
\begin{proof}
By Lemma~\ref{lem:fwk}, for every $\bsk\in \mathcal{E}(\bsell)$
    $$|\hat{f}^w(\bsk)|^2\leq 2^{|\supp(\bsell)|}\sum_{v\subseteq\supp(\bsell)}|\hat{f}^w_v(\bsk_v)|^2\prod_{j\in \supp(\bsell)\setminus v} |\hat{w}_j(k_j)|^2.$$
    Therefore,
    \begin{align*}
        &\sum_{\bsk\in \mathcal{E}(\bsell) } |\hat{f}^w(\bsk)|^2\prod_{j\in\supp(\bsell)} |\lambda(k_j)|^{-2\alpha}\\
        \leq & \sum_{\bsk\in \mathcal{E}(\bsell) } 2^{|\supp(\bsell)|}\sum_{v\subseteq\supp(\bsell)}|\hat{f}^w_v(\bsk_v)|^2\prod_{j\in v} |\lambda(k_j)|^{-2\alpha} \prod_{j\in \supp(\bsell)\setminus v} |\hat{w}_j(k_j)|^2 |\lambda(k_j)|^{-2\alpha} \\
        = & 2^{|\supp(\bsell)|}\sum_{v\subseteq\supp(\bsell)} \Big(\sum_{\bsk_v\in \mathcal{N}^{-1}_{|v|}(\bsell_v)}|\hat{f}^w_v(\bsk_v)|^2\prod_{j\in v} |\lambda(k_j)|^{-2\alpha} \Big) \prod_{j\in \supp(\bsell)\setminus v} \tilde{\mathcal{S}}(w_j,\ell_j,\alpha),
    \end{align*}
    where
    $$\tilde{\mathcal{S}}(w,\ell,\alpha)=\sum_{k\in \mathcal{N}^{-1}(\ell)} |\hat{w}(k)|^2 |\lambda(k)|^{-2\alpha}.$$
    By Lemma~\ref{lem:fkbound} and Hölder's inequality,
    \begin{align*}
        &\sum_{\bsk_v\in \mathcal{N}^{-1}_{|v|}(\bsell_v)}|\hat{f}^w_v(\bsk_v)|^2\prod_{j\in v} |\lambda(k_j)|^{-2\alpha}\\
        \leq &\Big(\sum_{\bsk_v\in \mathcal{N}^{-1}_{|v|}(\bsell_v)}|\hat{f}^w_v(\bsk_v)|^2\prod_{j\in v} |\lambda(k_j)|^{-2}\Big)^\alpha\Big(\sum_{\bsk_v\in \mathcal{N}^{-1}_{|v|}(\bsell_v)}|\hat{f}^w_v(\bsk_v)|^2\Big)^{1-\alpha}\\
        \leq & \Vert \partial^{v}f^w_v \Vert^{2\alpha}_{L^2(\itr^{|v|})}\Vert f^w_v \Vert^{2-2\alpha}_{L^2(\itr^{|v|})}.
    \end{align*}
    Meanwhile, by Lemma~\ref{lem:wkbound} and Equation~\eqref{eqn:intphik},
    \begin{align*}
       \tilde{\mathcal{S}}(w_\theta,\ell,\alpha)\leq  &16 \sum_{k\in \mathcal{N}^{-1}(\ell)}\min\big(\theta^2\Vert\phi_k\Vert^2_{L^\infty(\itr)}, \Vert \mathcal{I}(\phi_k)\Vert^2_{L^\infty(\itr)}\big)|\lambda(k)|^{-2\alpha} \\
       \leq &16C^2_\phi\sum_{k\in\mathcal{N}^{-1}(\ell) } \min\big(\theta^2 , |\lambda(k)|^{2}\big)|\lambda(k)|^{-2\alpha}. 
    \end{align*}
    The conclusion follows after we plug in the above bounds.
\end{proof}


\section{Norms of $f$ and $f^w$}\label{sec:norm}
In this section, we aim to characterize the norms of $f^w$ under various assumptions on $f$. We first consider the case $ f\in W^{1,q}_{\mix}(\real^s,\varphi)$.



\begin{lemma}\label{lem:dfw}
    For $f\in W^{1,q}_{\mix}(\real^s,\varphi)$ and $f^w(\bsu)=w(\bsu)f\circ T(\bsu)$,
    $$\partial^{1{:}s} f^w (\bsu) =\sum_{v\subseteq 1{:}s}(\partial^{v} f) \circ T(\bsu) \prod_{j\in v} \frac{w^2_j(u_j)}{\varphi\circ T_j(u_j)}\prod_{j\in v^c} w'_j(u_j).$$
\end{lemma}
\begin{proof}
 This is a direct application of chain rule for weak derivatives.
\end{proof}

Recall that $w_j=w_{\theta_j}$ is given by Equation~\eqref{eqn:wtheta} with $\theta_j\in (0,1/2]$ and $\eta=\teta_p$  given by \eqref{eq:tildeetap}, and $T_j$ is given by Equation~\eqref{eqn:Tj}. We summarize useful facts about $\eta_p$ in the following lemma.

\begin{lemma}\label{eqn:etabound}
For $p\geq 1$, $\eta_p(u)$ given by \eqref{eq:tildeetap} is monotonic, differentiable over $u\in [0,1/2]$ and satisfies
$$\eta_p(u)\geq \frac{u^{p+1}}{p} \Tilde{\eta}'_p(u) \quad\text{and} \quad \int_{0}^u\eta_p(t)\rd t\geq \frac{u^{p+1}}{p} \eta_p(u). $$
\end{lemma}
\begin{proof}
    See the appendix.
\end{proof}

\begin{lemma}\label{lem:rhoTbound}
    Suppose $\varphi$ satisfies Assumption~\ref{assump:rho}. For $\varepsilon>0$ and $u\in (0,1/2]$,
    $$\varphi\circ T_j(u)\geq 
       c_{\varepsilon,p} \Big(  \theta_jw_{j}(u)\min\big((2u/\theta_j)^{p+1},1 \big)\Big)^{1+\varepsilon},$$
    where $c_{\varepsilon,p}$ are constants depending on $\varepsilon$ and $p$.
\end{lemma}
\begin{proof}
    By Assumption~\ref{assump:rho} and Equation~\eqref{eqn:Tj}, for any $\varepsilon>0$,
    \begin{equation}\label{eqn:phiTjbound}
        \varphi\circ T_j(u)\geq c_\varepsilon\Big(\Phi\circ\Phi^{-1}\big(\mathcal{I}(w_{j})(u)\big) \Big)^{1+\epsilon}=c_\varepsilon\Big(\mathcal{I}(w_{j})(u)\Big)^{1+\epsilon}.
    \end{equation}
    When $u\in  (0,\theta_j/2] $, 
    $$\mathcal{I}(w_{j})(u)=(1-\theta_j)^{-1}\int_0^{u}\eta_p(t/\theta_j)\rd t=(1-\theta_j)^{-1}\theta_j \int_0^{u/\theta_j}\eta_p(t)\rd t.$$
     Lemma~\ref{eqn:etabound} implies that if $u\in  (0,\theta_j/2] $,
    $$\mathcal{I}(w_{j})(u)\geq
         \theta^{-p}_j u^{p+1}w_{j}(u)/p.$$
    When $u\in  (\theta_j/2,\theta_j) $, since $w_{j}(\theta_j/2)=w_{j}(\theta_j)/2\geq w_{j}(u)/2$,
    $$\mathcal{I}(w_{j})(u)\geq \mathcal{I}(w_{j})(\theta_j/2)\geq
         (\theta_j/2^{p+2}) w_{j}(u)/p.$$
    Finally when $u\in [\theta_j,1/2]$, $\mathcal{I}(w_{j})(u)=(1-\theta_j)^{-1}(u-\theta_j/2)$ and $w_{j}(u)=(1-\theta_j)^{-1}$, so $\mathcal{I}(w_{j})(u)\geq (\theta_j/2) w_{j}(u) $. Our conclusion follows after putting the above bounds into \eqref{eqn:phiTjbound}.
\end{proof}

\begin{lemma}\label{lem:dwjbound}
    For $u\in (0,1/2]$,
    $$|w'_j(u)|\leq 
      c_p  \theta^{-1}_j w_j(u) \max\big( (2u/\theta_j)^{-p-1},1\big),$$
    where $c_{p}$ are constants depending on $p$.
\end{lemma}
\begin{proof}
    Differentiating Equation~\eqref{eqn:wtheta} gives
    $$w'_{j}(u)=\begin{cases}
    \theta^{-1}_j(1-\theta_j)^{-1}\eta'_p(u/\theta_j), &\text{ if } u\in (0,\theta_j/2] \\
    \theta^{-1}_j(1-\theta_j)^{-1}\eta'_p(1-u/\theta_j) , &\text{ if } u\in (\theta_j/2,\theta_j) \\
    0, &\text{ if } u\in [\theta_j,1/2] 
\end{cases}.  $$
When $u\in (0,\theta_j/2]$, Lemma~\ref{eqn:etabound} implies
$$|w'_j(u)|\leq 
     p  \theta^{p}_j u^{-p-1} w_j(u).$$
When $u\in (\theta_j/2,\theta_j)$, 
$$|w'_j(u)|\leq \theta^{-1}_j(1-\theta_j)^{-1}\sup_{t\in [0,1/2]}\eta_p'(t)\leq 2\theta^{-1}_j w_j(u)\sup_{t\in [0,1/2]}\eta_p'(t).$$
The conclusion then follows since $\sup_{t\in [0,1/2]}\eta'_p(t)$ is finite and only depends on $p$. 
\end{proof}

\begin{theorem}\label{thm:Wqcase}
    For $f\in W^{1,q}_{\mix}(\real^s,\varphi)$ with $q>1$, $\varphi$ satisfying Assumption~\ref{assump:rho} and $\varepsilon\in (0,1-1/q)$, we have
    $$\Vert\partial^{1{:}s} f^w \Vert_{L^q(\itr^s)}\leq  \Big(\prod_{j=1}^sC_{\varepsilon,p,q}\theta^{-1-\varepsilon}_j\Big)\Vert f \Vert_{W^{1,q}_{\mix}(\real^s,\varphi)},$$
    where $C_{\varepsilon,p,q}$ is a constant depending on $\varepsilon,p$ and $q$.
\end{theorem}
\begin{proof}
    First we use Lemma~\ref{lem:dfw} to bound
    \begin{align}\label{eqn:partialfbound}
        |\partial^{1{:}s} f^w (\bsu)|\leq \sum_{v\subseteq 1{:}s}|(\partial^{v} f) \circ T(\bsu)| \prod_{j=1}^s w_j(u_j)^{1/q}\prod_{j\in v} \frac{w_j(u_j)^{2-1/q}}{\varphi\circ T_j(u_j)}\prod_{j\in v^c} \frac{|w'_j(u_j)|}{w_j(u_j)^{1/q}}.
    \end{align}
    By Lemma~\ref{lem:rhoTbound} and symmetry,
    \begin{align*}
        \sup_{u\in \itr }\frac{w_j(u)^{2-1/q}}{\varphi\circ T_j(u)}=&\sup_{u\in (0,1/2] }\frac{w_j(u)^{2-1/q}}{\varphi\circ T_j(u)}\\
        \leq &
       \tc^{-1}_{\varepsilon,p}  \theta^{-1-\varepsilon}_j\sup_{u\in (0,1/2] }w_{j}(u)^{1-1/q-\varepsilon}\max\big((2u/\theta_j)^{-(1+\varepsilon)(1+p)},1\big).
    \end{align*}
   Note that
   \begin{align*}
       &\sup_{u\in (0,1/2] }w_{j}(u)^{1-1/q-\varepsilon}\max\big((2u/\theta_j)^{-(1+\varepsilon)(1+p)},1\big)\\
       \leq & \max\Big(\sup_{u\in (0,\theta_j/2)}(2u/\theta_j)^{-(2-1/q)(1+p)}\exp\Big((1-1/q-\varepsilon)(2^p-(u/\theta_j)^{-p})\Big), 2\Big),
   \end{align*}
   which can be further bounded in terms of $\varepsilon,p$ and $q$. Hence
   \begin{equation}\label{eqn:woverphi}
      \sup_{u\in \itr }\frac{w_j(u)^{2-1/q}}{\varphi\circ T_j(u)}\leq C_1\theta^{-1-\varepsilon}_j 
   \end{equation}
   for $C_1$ depending on $\varepsilon,p$ and $q$. A similar calculation using Lemma~\ref{lem:dwjbound} shows
   \begin{equation}\label{eqn:dw}
       \sup_{u\in \itr } \frac{|w'_j(u_j)|}{w_j(u_j)^{1/q}} \leq C_2 \theta^{-1}_j
   \end{equation}
   for $C_2$ depending on $p$ and $q$. Using the above bounds, Equation~\eqref{eqn:partialfbound} becomes
   \begin{align*}
        |\partial^{1{:}s} f^w (\bsu)|\leq &\sum_{v\subseteq 1{:}s}|(\partial^{v} f) \circ T(\bsu)| \prod_{j=1}^s w_j(u_j)^{1/q} \prod_{j\in v} C_1\theta^{-1-\varepsilon}_j\prod_{j\in v^c} C_2\theta^{-1}_j\\
        \leq &  \Big(\prod_{j=1}^sC_{\varepsilon,p,q}\theta^{-1-\varepsilon}_j\Big) \Big(\sum_{v\subseteq 1{:}s}|(\partial^{v} f) \circ T(\bsu)|^q \prod_{j=1}^s w_j(u_j)\Big)^{1/q},
    \end{align*}
    where
    $$C_{\varepsilon,p,q} =\Big(\sum_{v\subseteq 1{:}s} C^{\frac{q}{q-1}|v|}_1C^{\frac{q}{q-1}(s-|v|)}_2\Big)^{\frac{1}{s}(1-\frac{1}{q})}=\Big(C^{\frac{q}{q-1}}_1+C^{\frac{q}{q-1}}_2\Big)^{1-\frac{1}{q}}.$$
    The conclusion then follows because
    \begin{align*}
     &\sum_{v\subseteq 1{:}s}\int_{\itr^s}|(\partial^{v} f) \circ T(\bsu)|^q \prod_{j=1}^s w_j(u_j)\rd \bsu\\
     =&\sum_{v\subseteq 1{:}s}\int_{\itr^s}|(\partial^{v} f) \circ T(\bsu)|^q \prod_{j=1}^s \varphi\circ T_j(u_j)T'_j(u_j)\rd \bsu\\
     = &\sum_{v\subseteq 1{:}s}\int_{\real^s}|\partial^{v} f(\bsx)|^q \prod_{j=1}^s \varphi(x_j)\rd \bsx=\Vert f \Vert^q_{W^{1,q}_{\mix}(\real^s,\varphi)}.
    \end{align*}
\end{proof}

Theorem~\ref{thm:Wqcase} shows $\partial^{1{:}s} f^w \in L^2(\itr^s)$ when $f\in W^{1,2}_{\mix}(\real^s,\varphi)$. However, in some applications the boundary growth of $f$ is too rapid for even $f\in L^2(\real^s,\varphi)$ to hold. Characterizing $f$ by its $W^{1,q,\infty}_{\mix}(\real^s,\varphi)$-norm is a better choice for such cases. The next theorem shows $f^w \in W^{1,2}_{\mix}(\itr^s)$ even for $f\in W^{1,q,\infty}_{\mix}(\real^s,\varphi)$ with $1<q\leq 2$.

\begin{theorem}\label{thm:Wqinfcase}
     For $f\in W^{1,q,\infty}_{\mix}(\real^s,\varphi)$ with $1<q\le 2$, $\varphi$ satisfying Assumption~\ref{assump:rho}, $\varepsilon\in (0,(q-1)/(q+1))$ and $q'\geq q$, then
     \begin{equation}\label{eqn:fLq'bound}
       \Vert f^w \Vert_{L^{q'}(\itr^s)} \leq  \Big(\prod_{j=1}^sC_{\varepsilon,p,q,q'}\theta^{-(1+\varepsilon)(1/q-1/q')}_j\Big)\Vert f \Vert_{L^{q,\infty}(\real^s,\varphi)}   
     \end{equation}
     and
     \begin{equation}\label{eqn:DfLq'bound}
         \Vert\partial^{1{:}s} f^w \Vert_{L^{q'}(\itr^s)}\leq  \Big(\prod_{j=1}^sC'_{\varepsilon,p,q,q'}\theta^{-(1+\varepsilon)(1+1/q-1/q')}_j\Big)\Vert f \Vert_{W^{1,q,\infty}_{\mix}(\real^s,\varphi)},
     \end{equation}
    where $C_{\varepsilon,p,q,q'},C'_{\varepsilon,p,q,q'}$ are constants depending on $\varepsilon,p,q$ and $q'$.
\end{theorem}
\begin{proof}
We first prove the bound on $ \Vert f^w \Vert_{L^{q'}(\itr^s)}  $. Let $\varepsilon'\in (0,q-1)$. By Lemma~\ref{lem:embedding},  $\Vert f \Vert_{L^{q-\varepsilon'}(\itr^s)}\leq C^s_3\Vert f \Vert_{L^{q,\infty}(\itr^s)}$ for $C_3$ depending on $q$ and $\varepsilon'$. By Equation~\eqref{eqn:fwLqbound}, 
    $$\Vert f^w \Vert_{L^{q-\varepsilon'}(\itr^s)} \leq 2^{\frac{q-\varepsilon'-1}{q-\varepsilon'}s}\Vert f \Vert_{L^{q-\varepsilon'}(\real^s,\varphi)}\leq \Big(2^{\frac{q-\varepsilon'-1}{q-\varepsilon'}}C_3\Big)^s\Vert f \Vert_{L^{q,\infty}(\real^s,\varphi)}.$$
    Next, we use the definition of $L^{q,\infty}(\real^s,\varphi)$-norm to bound
    $$|f^w (\bsu)| = |f\circ T(\bsu)|\prod_{j=1}^s w_j(u_j) \leq \varphi^{s/q}_\infty\Vert f \Vert_{L^{q,\infty}(\real^s,\varphi)}\prod_{j=1}^s \frac{w_j(u_j)}{\varphi\circ T_j(u_j)^{1/q}}.  $$
    By Equation~\eqref{eqn:woverphi} with $q^*=1/(2-q)$,
    $$\sup_{u\in \itr}\frac{w_j(u)}{\varphi\circ T_j(u)^{1/q}}=\Big(\sup_{u\in \itr}\frac{w_j(u)^{2-1/q^*}}{\varphi\circ T_j(u)}\Big)^{1/q} \leq C_4 \theta_j^{-(1+\varepsilon/2)/q}$$
for $C_4$ depending on $\varepsilon,p$ and $q$. Hence,
$$\Vert f^w \Vert_{L^{\infty}(\itr^s)} \leq \varphi^{s/q}_\infty\Vert f \Vert_{L^{q,\infty}(\real^s,\varphi)} \prod_{j=1}^s C_4\theta_j^{-(1+\varepsilon/2)/q}. $$
Then by the interpolation inequality \cite[Theorem 2.11]{adams2003sobolev},
$$\Vert f^w \Vert_{L^{q'}(\itr^s)} \leq \Vert f^w \Vert_{L^{q-\varepsilon'}(\itr^s)}^{\frac{q-\varepsilon'}{q'}} \Vert f^w \Vert^{1-\frac{q-\varepsilon'}{q'}}_{L^{\infty}(\itr^s)}\leq \Vert f \Vert_{L^{q,\infty}(\real^s,\varphi)}\prod_{j=1}^s C_5 \theta_j^{-A}$$
for $C_5$ depending on $\varepsilon,\varepsilon',p,q$ and $q'$ (note that $\varphi_\infty$ is treated as a constant), and
$$A=\Big(1-\frac{q-\varepsilon'}{q'}\Big)\frac{1+\varepsilon/2}{q}=(1+\varepsilon)\Big(\frac{1}{q}-\frac{1}{q'}\Big)-\frac{\varepsilon}{2}\Big(\frac{1}{q}-\frac{1}{q'}-\frac{2\varepsilon'}{\varepsilon qq'}-\frac{\varepsilon'}{ qq'}\Big).$$
Equation~\eqref{eqn:fLq'bound} follows after we choose a sufficiently small $\varepsilon'$.

The bound on $\Vert\partial^{1{:}s} f^w \Vert_{L^{q'}(\itr^s)}$ can be proven similarly. By Lemma~\ref{lem:embedding} and Theorem~\ref{thm:Wqcase},
 $$\Vert\partial^{1{:}s} f^w \Vert_{L^{q-\varepsilon'}(\itr^s)} \leq  \Big(\prod_{j=1}^sC_6\theta^{-1-\varepsilon}_j\Big)\Vert f \Vert_{W^{1,q,\infty}_{\mix}(\real^s,\varphi)}$$
    for $C_6$ depending on $\varepsilon,\varepsilon',p$ and $q$.
    Next, we use Lemma~\ref{lem:dfw} to bound
    \begin{align*}
        |\partial^{1{:}s} f^w (\bsu)|\leq & \sum_{v\subseteq 1{:}s}\Vert \partial^v f \Vert_{L^{q,\infty}(\real^s,\varphi)}\Big(\prod_{j=1}^s\frac{\varphi \circ T_j(\bsu)}{\varphi_\infty}\Big)^{-1/q} \prod_{j\in v} \frac{w_j(u_j)^2}{\varphi\circ T_j(u_j)}\prod_{j\in v^c} |w'_j(u_j)|\\
        =& \varphi^{s/q}_\infty\sum_{v\subseteq 1{:}s}\Vert \partial^v f \Vert_{L^{q,\infty}(\real^s,\varphi)} \prod_{j\in v} \frac{w_j(u_j)^2}{(\varphi\circ T_j(u_j))^{1+1/q}}\prod_{j\in v^c} \frac{|w'_j(u_j)|}{(\varphi\circ T_j(u_j))^{1/q}}.
    \end{align*}
    By Equation~\eqref{eqn:woverphi} with $q^*$ satisfying $(2-1/q^*)(1+1/q)=2$ and $\varepsilon\in (0,1-1/q^*)$,
    $$\sup_{u\in \itr}\frac{w_j(u)^2}{(\varphi\circ T_j(u))^{1+1/q}}=\Big(\sup_{u\in \itr} \frac{w_j(u)^{2-1/q^*}}{\varphi\circ T_j(u)}\Big)^{1+1/q}\leq C_7 \theta_j^{-(1+\varepsilon/2)(1+1/q)}$$
    for $C_7$ depending on $\varepsilon,p$ and $q$, where $p$ is required to satisfy $p\geq (1-1/q^*-\varepsilon)^{-1}(1+\varepsilon)$ if $\eta=\eta_p$.
     Similarly by Equation~\eqref{eqn:dw}, 
     $$\sup_{u\in \itr}\frac{|w'_j(u_j)|}{(\varphi\circ T_j(u_j))^{1/q}}\leq \sup_{u\in \itr}\frac{|w'_j(u_j)|}{w_j(u_j)^{2/(q+1)}} \Big(\sup_{u\in \itr} \frac{w_j(u)^{2-1/q^*}}{\varphi\circ T_j(u)}\Big)^{1/q}\leq C_8 \theta_j^{-1-(1+\varepsilon/2)/q}$$
     for $C_8$ depending on $\varepsilon,p$ and $q$. Using the above bounds,
     \begin{align*}
       & \Vert\partial^{1{:}s} f^w\Vert_{L^{\infty}(\itr)}\\
       \leq & \varphi^{s/q}_\infty\sum_{v\subseteq 1{:}s}\Vert \partial^v f \Vert_{L^{q,\infty}(\real^s,\varphi)}  \prod_{j\in v}C_7\theta_j^{-(1+\varepsilon/2)(1+1/q)} \prod_{j\in v^c}C_8 \theta_j^{-1-(1+\varepsilon/2)/q}\\
        \leq & \varphi^{s/q}_\infty\Vert f \Vert_{W^{1,q,\infty}_{\mix}(\real^s,\varphi)}\Big(\sum_{v\subseteq 1{:}s} C^{\frac{q}{q-1}|v|}_7C^{\frac{q}{q-1}(s-|v|)}_8\Big)^{1-\frac{1}{q}} \prod_{j=1}^s \theta_j^{-(1+\varepsilon/2)(1+1/q)}\\
        = & \varphi^{s/q}_\infty\Vert f \Vert_{W^{1,q,\infty}_{\mix}(\real^s,\varphi)} \prod_{j=1}^s \Big(C^{\frac{q}{q-1}}_7+C^{\frac{q}{q-1}}_8\Big)^{1-\frac{1}{q}} \theta_j^{-(1+\varepsilon/2)(1+1/q)}.
    \end{align*}
    Finally, we use the interpolation inequality to get 
    \begin{align*}
        \Vert\partial^{1{:}s} f^w \Vert_{L^{q'}(\itr^s)}\leq &\Vert\partial^{1{:}s} f^w \Vert_{L^{q-\varepsilon'}(\itr^s)}^{\frac{q-\varepsilon'}{q'}} \Vert\partial^{1{:}s} f^w \Vert^{1-\frac{q-\varepsilon'}{q'}}_{L^{\infty}(\itr^s)}\\
        \leq &\Vert f \Vert_{W^{1,q,\infty}_{\mix}(\real^s,\varphi)} \prod_{j=1}^s C_9 \theta_j^{-B}
    \end{align*}
    for $C_9$ depending on $\varepsilon,\varepsilon',p,q$ and $q'$, and
    \begin{align*}
        B=&\frac{q-\varepsilon'}{q'}\Big(1+\varepsilon\Big)+\Big(1-\frac{q-\varepsilon'}{q'}\Big)\Big(1+\frac{1}{q}\Big)\Big(1+\frac{\varepsilon}{2}\Big)\\
        =&\Big(1+\frac{1}{q}-\frac{1}{q'}\Big)\Big(1+\varepsilon\Big)-\frac{\varepsilon}{2qq'}\Big((q'-q)(1+q)-\frac{\varepsilon'}{\varepsilon}(2+\varepsilon-\varepsilon q)\Big).
    \end{align*}
    Equation~\eqref{eqn:DfLq'bound} follows after we choose a sufficiently small $\varepsilon'$.

     
\end{proof}



    %%\section{Application to lattice rules}

%%In this section, we apply our theory and derive conditions under which $f^w$ can be efficiently integrated by lattice rules. A key step is to show $f^w$ lies in a Korobov space with modified weights.

%%Let $\{\phi_k\mid k\in \natu_0\}$ be the orthonormal basis of trigonometric functions. One way is to set $\overline{\phi_k(u)}=\cos(\pi k u)$ for even $k$ and $\overline{\phi_k(u)}=\sin(\pi (k+1)u)$ for odd $k$. However, it is more convenient for our proof to change the index from $\natu_0$ to integers $\ints$ and set $\overline{\phi_k(u)}=\exp(2\pi\icomp k u)$ with $\icomp$ the imaginary number $\ii$. While we cannot directly apply Lemma~\ref{lem:fkbound} to trigonometric functions, the following modified version applies:

\section{Application to digital nets}\label{sec:main}

In this section, we apply our theory and derive conditions under which $f^w$ can be efficiently integrated by digital nets. In Example~\ref{example}, we have shown that $\phi_k(u)={}_b\walk (u)$ satisfies Equation~\eqref{eqn:intphik} for $\mathcal{N}(k)=\kappa_r b^{r-1}$ and $\lambda(k)=\Vert\mathcal{I}(\phi_{\mathcal{N}(k)})\Vert_{L^\infty(\itr)}\leq b^{-r+1}$, where $r$ and $\kappa_r$ are determined by $\kappa_r b^{r-1}\leq k <(\kappa_r+1) b^r$. Combining Theorem~\ref{thm:fkbound} and Theorem~\ref{thm:Wqcase}, we arrive at the following theorem.
\begin{theorem}\label{thm:Lqfwkbound}
For $\phi_k(u)={}_b\walk (u)$, $\alpha\in (0,1/2)$, $\bsell\in \natu_0^{s}\setminus\{\bszero\}$, $\bskappa\in \{1,\dots,b-1\}^{s}$, $f\in W^{1,2}_{\mix}(\real^s,\varphi)$, $w_j=w_{\theta_j}$  for $\{\theta_j\mid j\in 1{:}s\}\subseteq (0,1/2]$ and $\varepsilon\in(0,1/2)$, we have
    \begin{multline*}
        \sum_{\bsk\in \mathcal{E}(\bsell,\bskappa) } |\hat{f}^w(\bsk)|^2
        \leq \\
        C_{\varepsilon,p,b,\alpha}^{|\supp(\bsell)|} b^{-2\alpha\Vert \bsell\Vert_1}\sum_{v\subseteq\supp(\bsell)}\Vert f_v \Vert^2_{W^{1,2}_{\mix}(\real^s,\varphi)}\prod_{j\in v} \theta^{-2\alpha(1+\varepsilon)}_j\prod_{j\in \supp(\bsell)\setminus v} \theta^{1-2\alpha}_j,
    \end{multline*}
    where $C_{\varepsilon,p,b,\alpha}$ is a constant depending on $\varepsilon,p,b$ and $\alpha$, and
    $$\mathcal{E}(\bsell,\bskappa)=\Big\{\bsk'\in \natu_0^{s}\mid \supp(\bsk')=\supp(\bsell), \kappa_j b^{\ell_j-1} \leq k'_j < (\kappa_j+1)b^{\ell_j-1} \ \forall j\in \supp(\bsell)\Big\}.$$
\end{theorem}

\begin{proof}
For $\phi_k(u)={}_b\walk (u)$, $C_\phi=\sup_{k\in \natu_0} \Vert \phi_k\Vert_{L^\infty(\itr)}=1$ and $\mathcal{N}^{-1}(\kappa_j b^{\ell_j-1})=\{k'\in\natu \mid \kappa_j b^{\ell_j-1} \leq k' < (\kappa_j+1)b^{\ell_j-1}\}$. It follows that $\lambda(k')=\lambda(\kappa_j b^{\ell_j-1})=\Vert\mathcal{I}(\phi_{\kappa_j b^{\ell_j-1}})\Vert_{L^\infty(\itr)}\leq b^{-\ell_j+1} $ for $k'\in \mathcal{N}^{-1}(\kappa_j b^{\ell_j-1})$. Hence
\begin{align*}
\mathcal{S}(\theta_j,\kappa_j b^{\ell_j-1},\alpha)=&16\sum_{k'\in\mathcal{N}^{-1}(\kappa_j b^{\ell_j-1}) } \min\big(\theta_j^2 , |\lambda(k')|^{2}\big)|\lambda(k')|^{-2\alpha} \\
=& 16 b^{\ell_j-1} \lambda(\kappa_j b^{\ell_j-1})^{-2\alpha} \min\big(\theta_j^2, \lambda(\kappa_j b^{\ell_j-1})^{2}\big) \\
\leq & 16  \min\big(\theta_j^2\lambda(\kappa_j b^{\ell_j-1})^{-1-2\alpha}, \lambda(\kappa_j b^{\ell_j-1})^{1-2\alpha}\big)\\
\leq & 16 \theta^{1-2\alpha}_j,
\end{align*}
where we have used $1-2\alpha>0$. Next, by Equation~\eqref{eqn:fwLqbound} with $q=2$,
    \begin{align}\label{eqn:fL2bound}
        \Vert f^w_v \Vert_{L^2(\itr^{|v|})}
        \leq   2^{|v|/2} \Vert f_v \Vert_{L^2(\real^s,\varphi)} \leq 2^{|v|/2} \Vert f_v \Vert_{W^{1,2}_{\mix}(\real^s,\varphi)}.
    \end{align}
    Meanwhile, by Theorem~\ref{thm:Wqcase} with $f=f_v$, $q=2$ and identifying $v$ with $1{:}|v|$,
    \begin{equation}\label{eqn:DfL2bound}
        \Vert\partial^{v} f_v^w \Vert_{L^2(\itr^{|v|})}\leq  \Big(\prod_{j\in v}C_{\varepsilon,p,2}\theta^{-1-\varepsilon}_j\Big)\Vert f_v \Vert_{W^{1,2}_{\mix}(\real^s,\varphi)}.
    \end{equation}
    Using the above bounds, Equation~\eqref{eqn:fwkbound} becomes
    \begin{align*}
        &\sum_{\bsk\in \mathcal{E}(\bsell,\bskappa) } |\hat{f}^w(\bsk)|^2\prod_{j\in\supp(\bsell)} |\lambda(k_j)|^{-2\alpha} =\Big(\prod_{j\in\supp(\bsell)} \lambda(\kappa_j b^{\ell_j-1}) \Big)^{-2\alpha} \sum_{\bsk\in \mathcal{E}(\bsell,\bskappa) } |\hat{f}^w(\bsk)|^2 \\
      \leq &  2^{|\supp(\bsell)|}\sum_{v\subseteq\supp(\bsell)}2^{(1-\alpha)|v|}\Big(\prod_{j\in v}C_{\varepsilon,p,2}\theta^{-1-\varepsilon}_j\Big)^{2\alpha}\Vert f_v \Vert^2_{W^{1,2}_{\mix}(\real^s,\varphi)} \prod_{j\in \supp(\bsell)\setminus v} 16 \theta^{1-2\alpha}_j\\
      \leq & \max\big( 2^{2-\alpha} C^{2\alpha}_{\varepsilon,p,2} , 32\big)^{|\supp(\bsell)|} \sum_{v\subseteq\supp(\bsell)}\Vert f_v \Vert^2_{W^{1,2}_{\mix}(\real^s,\varphi)}\prod_{j\in v} \theta^{-2\alpha(1+\varepsilon)}_j\prod_{j\in \supp(\bsell)\setminus v} \theta^{1-2\alpha}_j.
    \end{align*}
    After bounding $\lambda(\kappa_j b^{\ell_j-1}) \leq b^{-\ell_j+1}$, we conclude
    \begin{multline*}
     \sum_{\bsk\in \mathcal{E}(\bsell,\bskappa) } |\hat{f}^w(\bsk)|^2 \leq \\ C_{\varepsilon,p,b,\alpha}^{|\supp(\bsell)|} b^{-2\alpha\Vert\bsell\Vert_1}  \sum_{v\subseteq\supp(\bsell)}\Vert f_v \Vert^2_{W^{1,2}_{\mix}(\real^s,\varphi)}\prod_{j\in v} \theta^{-2\alpha(1+\varepsilon)}_j\prod_{j\in \supp(\bsell)\setminus v} \theta^{1-2\alpha}_j
    \end{multline*}
    with 
    $C_{\varepsilon,p,b,\alpha}=\max\big( 2^{2-\alpha} C^{2\alpha}_{\varepsilon,p,2}  , 32\big)b^{2\alpha}$.
\end{proof}

 %Finally, because $u_{r}$ in $u=\sum_{i=1}^\infty u_i b^{-i}$ equals $0$ over $u\in [0,b^{-r})$,
    %$$\lambda(\kappa_j b^{\ell_j-1})=\Vert\mathcal{I}(\phi_{\kappa_j b^{\ell_j-1}})\Vert_{L^\infty(\itr)}\geq \Big|\int_0^{b^{-\ell_j}}  \exp\Big(\frac{2\pi \ii}{b} \kappa_j u_{\ell_j}\Big) \rd u\Big|=b^{-\ell_j}.$$

The next theorem is the counterpart of Theorem~\ref{thm:Lqfwkbound} for $f\in W^{1,q,\infty}_{\mix}(\real^s,\varphi)$.

\begin{theorem}\label{thm:Lqinffwkbound}
For $\phi_k(u)={}_b\walk (u)$, $\alpha\in (0,1/2)$, $\bsell\in \natu_0^{s}\setminus\{\bszero\}$, $\bskappa\in \{1,\dots,b-1\}^{s}$, $f\in W^{1,q,\infty}_{\mix}(\real^s,\varphi)$ with $q\in (1,2]$, $w_j=w_{\theta_j}$ for  $\{\theta_j\mid j\in 1{:}s\}\subseteq (0,1/2]$ and $\varepsilon\in (0,(q-1)/(q+1))$, we have
    \begin{multline*}
        \sum_{\bsk\in \mathcal{E}(\bsell,\bskappa) } |\hat{f}^w(\bsk)|^2
        \leq \\
        C_{\varepsilon,p,q,b,\alpha}^{|\supp(\bsell)|} b^{-2\alpha\Vert \bsell\Vert_1}\sum_{v\subseteq\supp(\bsell)}\Vert f_v \Vert^2_{W^{1,q,\infty}_{\mix}(\real^s,\varphi)}\prod_{j\in v} \theta^{-(1+\varepsilon)(2\alpha-1+2/q)}_j\prod_{j\in \supp(\bsell)\setminus v} \theta^{1-2\alpha}_j,
    \end{multline*}
    where $C_{\varepsilon,p,q,b,\alpha}$ is a constant depending on $\varepsilon,p,q,b$ and $\alpha$.
\end{theorem}

    \begin{proof}
        The proof is essentially the same as that of Theorem~\ref{thm:Lqfwkbound}, except we replace Equation~\eqref{eqn:fL2bound} with
        $$\Vert f_v^w \Vert_{L^{2}(\itr^{|v|})} \leq  \Big(\prod_{j\in v}C_{\varepsilon,p,q,2}\theta^{-(1+\varepsilon)(1/q-1/2)}_j\Big)\Vert f_v \Vert_{W^{1,q,\infty}_{\mix}(\real^s,\varphi)}  $$
    and Equation~\eqref{eqn:DfL2bound} with
    $$\Vert\partial^{v} f_v^w \Vert_{L^{2}(\itr^{|v|})}\leq  \Big(\prod_{j\in v}C'_{\varepsilon,p,q,2}\theta^{-(1+\varepsilon)(1/q+1/2)}_j\Big)\Vert f_v \Vert_{W^{1,q,\infty}_{\mix}(\real^s,\varphi)}, $$
    where $C_{\varepsilon,p,q,2}$ and $C'_{\varepsilon,p,q,2}$ come from Theorem~\ref{thm:Wqinfcase}.
    \end{proof}



    We are ready to bound the variance of $\hat{\mu}$ for digital nets. To simplify the notation, we let 
    \begin{equation*}
     \Vert f\Vert_{q}=\begin{cases}
        \sup_{ v\subseteq 1{:}s}\Vert f_v \Vert_{W^{1,q}_{\mix}(\real^s,\varphi)}, &\text{ if } f\in W^{1,q}_{\mix}(\real^s,\varphi)\\
        \sup_{v\subseteq 1{:}s}\Vert f_v \Vert_{W^{1,q,\infty}_{\mix}(\real^s,\varphi)},&\text{ if } f\in W^{1,q,\infty}_{\mix}(\real^s,\varphi)
    \end{cases},   
    \end{equation*}
    and for $v \subseteq 1{:}s$
    \begin{equation*}
     \gamma_v=\begin{cases}
        \Vert f\Vert^{-1}_{q} \Vert f_v \Vert_{W^{1,q}_{\mix}(\real^s,\varphi)}, &\text{ if } f\in W^{1,q}_{\mix}(\real^s,\varphi)\\
        \Vert f\Vert^{-1}_{q} \Vert f_v \Vert_{W^{1,q,\infty}_{\mix}(\real^s,\varphi)},&\text{ if } f\in W^{1,q,\infty}_{\mix}(\real^s,\varphi)
    \end{cases}.   
    \end{equation*}

    \begin{theorem}\label{thm:qmcvar}
    Let $\{\bsu_0,\dots,\bsu_{n-1}\}$ be a scrambled digital net in base $b\ge 2$ with $t$-quality parameters $\{t_\omega\mid \emptyset\neq \omega\subseteq 1{:}s\}$.
        If $f\in W^{1,q}_{\mix}(\real^s,\varphi)$ with $q=2$ and $w_j,\varepsilon$ satisfy the assumptions of Theorem~\ref{thm:Lqfwkbound}, or if $f\in W^{1,q,\infty}_{\mix}(\real^s,\varphi)$ with $q\in (1,2]$ and $w_j,\varepsilon$ satisfy the assumptions of Theorem~\ref{thm:Lqinffwkbound}, then for any $\alpha\in (0,1/2)$, 
        $$ \e \Big|\frac{1}{n}\sum_{i=0}^{n-1}f^w(\bsu_i)-\int_{\real^s} f(\bsx) \prod_{j=1}^s \varphi(x_j)\rd \bsx\Big|^2 
            \leq \frac{\Vert f\Vert^2_q }{b^{(1+2\alpha)m}}\sum_{\emptyset \neq \omega\subseteq 1{:}s }  C^{|\omega|}_{*} m^{|\omega|-1}\Tilde{\gamma}_\omega,$$
        where $C_{*}$ is a constant depending on $\varepsilon,p,q,b$ and $\alpha$, and
        $$\Tilde{\gamma}_\omega=b^{(1+2\alpha)t_w} \sum_{v\subseteq \omega}\gamma_v^2\prod_{j\in v} \theta^{-(1+\varepsilon)(2\alpha-1+2/q)}_j \prod_{j\in \omega\setminus v} \theta^{1-2\alpha}_j.$$
    \end{theorem}
    \begin{proof}
        First we notice that $L_{\bsell}$ defined by Equation~\eqref{eqn:Ll} is the union of $\mathcal{E}(\bsell,\bskappa)$ for $\bskappa$ ranging over $\{1,\dots,b-1\}^s$. Since $\mathcal{E}(\bsell,\bskappa)$ does not depend on $\kappa_j$ for $j\notin \supp(\bsell)$, there are $(b-1)^{|\supp(\bsell)|}$ number of disjoint $\mathcal{E}(\bsell,\bskappa)$ and 
        \begin{align*}
            \sigma^2_{\bsell}=&\sum_{\bsk\in L_{\bsell}} |\hat{f}(\bsk)|^2 \\
            \leq & C_{10}^{|\supp(\bsell)|} b^{-2\alpha\Vert \bsell\Vert_1}\Vert f\Vert^2_q\sum_{v\subseteq\supp(\bsell)} \gamma_v^2 \prod_{j\in v} \theta^{-(1+\varepsilon)(2\alpha-1+2/q)}_j\prod_{j\in \supp(\bsell)\setminus v} \theta^{1-2\alpha}_j,
        \end{align*}
        where $C_{10}=(b-1)C_{\varepsilon,p,b,\alpha}$ if $f\in W^{1,q}_{\mix}(\real^s,\varphi)$ and $C_{10}=(b-1)C_{\varepsilon,p,q,b,\alpha}$ if $f\in W^{1,q,\infty}_{\mix}(\real^s,\varphi)$. 
        Then by Lemma~\ref{lem:qmcvar}, 
        \begin{align*}
            &\e \Big|\frac{1}{n}\sum_{i=0}^{n-1}f^w(\bsu_i)-\int_{\itr^s} f^w(\bsu)\rd \bsu\Big|^2  \\
            \leq & \frac{1}{b^m} \sum_{\emptyset\neq \omega \subseteq 1{:}s}\sum_{\bsell\in \natu^\omega}\Gamma_{\omega,\bsell} C_{10}^{|\omega|} b^{-2\alpha\Vert \bsell\Vert_1}\Vert f\Vert^2_q\sum_{v\subseteq \omega}\gamma_v^2\prod_{j\in v} \theta^{-(1+\varepsilon)(2\alpha-1+2/q)}_j\prod_{j\in \omega\setminus v} \theta^{1-2\alpha}_j\nonumber\\
            =&\frac{\Vert f\Vert^2_q}{b^m} \sum_{\emptyset\neq \omega \subseteq 1{:}s}C_{10}^{|\omega|}\Big(\sum_{\bsell\in \natu^\omega} \Gamma_{\omega,\bsell}b^{-2\alpha\Vert \bsell\Vert_1}\Big)\sum_{v\subseteq \omega}\gamma_v^2\prod_{j\in v} \theta^{-(1+\varepsilon)(2\alpha-1+2/q)}_j\prod_{j\in \omega\setminus v} \theta^{1-2\alpha}_j.\nonumber
        \end{align*}
        By Equation~\eqref{eqn:gaincoefbound}, 
        \begin{align*}
          \sum_{\bsell\in \natu^\omega} \Gamma_{\omega,\bsell}b^{-2\alpha\Vert \bsell\Vert_1}\leq &\Big(\frac{b}{b-1}\Big)^{|\omega|-1}b^{t_{\omega}}\sum_{\bsell\in \natu^{\omega}}b^{-2\alpha\Vert \bsell\Vert_1}\bsone\{\Vert \bsell\Vert_1> m-t_{\omega}-|\omega|\}.
        \end{align*}
        Because there are ${N-1\choose |\omega|-1}$ number of $\bsell\in \natu^{\omega}$ satisfying $\Vert \bsell\Vert_1=N$ for $N\geq |\omega|$, 
        \begin{align*}
            &\sum_{\bsell\in \natu^{\omega}}b^{-2\alpha\Vert \bsell\Vert_1} \bsone\{\Vert \bsell\Vert_1> m-t_{\omega}-|\omega|\}\leq \sum_{N=\max(m-t_{\omega}-|\omega|,|\omega|)}^\infty {N-1\choose |\omega|-1}b^{-2\alpha N}\\
            \leq & \frac{b^{-2\alpha (m-t_\omega-|\omega|)}}{(1-b^{-2\alpha})^{|\omega|}}  \max\left({m-t_\omega-|\omega|-1\choose |\omega|-1},1\right),
        \end{align*}
        where the last inequality follows from \cite[Lemma 13.24]{dick:pill:2010}. Plugging the above bound,
        \begin{align*}
          \sum_{\bsell\in \natu^\omega} \Gamma_{\omega,\bsell}b^{-2\alpha\Vert \bsell\Vert_1}\leq &\Big(\frac{b}{b-1}\Big)^{|\omega|-1}b^{t_{\omega}}\frac{b^{-2\alpha (m-t_\omega-|\omega|)}}{(1-b^{-2\alpha})^{|\omega|}} m^{|\omega|-1}\leq \frac{C^{|\omega|}_{11}m^{|\omega|-1}b^{(1+2\alpha)t_\omega}}{b^{2\alpha m}}
        \end{align*}
        for $C_{11}$ depending on $b$ and $\alpha$, and
        \begin{align*}
            &\e \Big|\frac{1}{n}\sum_{i=0}^{n-1}f^w(\bsu_i)-\int_{\itr^s} f^w(\bsu)\rd \bsu\Big|^2  \\
            \leq 
            &\frac{\Vert f\Vert^2_q}{b^{(1+2\alpha)m}}\sum_{\emptyset\neq \omega \subseteq 1{:}s}C_{*}^{|\omega|}m^{|\omega|-1}b^{(1+2\alpha)t_\omega}\sum_{v\subseteq \omega}\gamma_v^2\prod_{j\in v} \theta^{-(1+\varepsilon)(2\alpha-1+2/q)}_j\prod_{j\in \omega\setminus v} \theta^{1-2\alpha}_j
        \end{align*}
        for $C_{*}=C_{10}C_{11}$. The conclusion follows from Equation~\eqref{eqn:intfw}.
    \end{proof}
\begin{remark}
    If $f\in W^{1,q,\infty}_{\mix}(\real^s,\varphi)$ with $q>2$, applying Lemma~\ref{lem:embedding} gives $f\in W^{1,q}_{\mix}(\real^s,\varphi)$ with $q=2$ so that Theorem~\ref{thm:qmcvar} applies. As shown in Remark~\ref{rem:rqmcrate}, when the density is standard Gaussian, simple RQMC without IS has a root mean squared error rate of $O(n^{-1+1/q+\epsilon})$ for arbitrarily small $\epsilon>0$, while our proposed boundary-damping IS improves the rate to $O(n^{-1+\epsilon})$.
\end{remark}

    \begin{corollary}\label{cor:qmcvar}
        Suppose $t$-quality parameters $\{t_\omega\mid \emptyset\neq \omega\subseteq 1{:}s\}$  satisfy Equation~\eqref{eqn:tutj} and $\{\gamma_v\mid v\subseteq 1{:}s\}$ satisfy 
        $$\gamma_v\leq \prod_{j\in v}\Gamma_j \ \forall v\subseteq 1{:}s$$
        for $\{\Gamma_j\mid j\in 1{:}s\}$. 
        Then under the assumptions of Theorem~\ref{thm:qmcvar},
        $$\e \Big|\frac{1}{n}\sum_{i=0}^{n-1}f^w(\bsu_i)-\int_{\real^s} f(\bsx) \prod_{j=1}^s \varphi(x_j)\rd \bsx\Big|^2 \leq \frac{\Vert f\Vert^2_q }{b^{(1+2\alpha)m}m}\prod_{j=1}^s \Big(1+C_*\Tilde{\Gamma}_j m\Big),$$
        where $C_*$ comes from Theorem~\ref{thm:qmcvar} and
        $$\Tilde{\Gamma}_j=b^{(1+2\alpha)t_j} (\Gamma^2_j \theta^{-(1+\epsilon)(2\alpha-1+2/q) }_j+\theta^{1-2\alpha}_j).$$
    \end{corollary}
    \begin{proof}
        First we compute
        \begin{align*}
            \Tilde{\gamma}_\omega\le&\Big(\prod_{j\in\omega }b^{(1+2\alpha)t_j}\Big) \sum_{v\subseteq \omega}\prod_{j\in v} \Gamma^2_j\theta^{-(1+\varepsilon)(2\alpha-1+2/q)}_j \prod_{j\in \omega\setminus v} \theta^{1-2\alpha}_j.\\
            =&\Big(\prod_{j\in\omega }b^{(1+2\alpha)t_j}\Big) \prod_{j\in \omega}\Big(\Gamma^2_j\theta^{-(1+\varepsilon)(2\alpha-1+2/q)}_j+\theta^{1-2\alpha}_j\Big)
            =\prod_{j\in\omega}\Tilde{\Gamma}_j.
        \end{align*}
        The conclusion follows from Theorem~\ref{thm:qmcvar} and
        \begin{align*}
         \sum_{\emptyset \neq \omega\subseteq 1{:}s }  C^{|\omega|}_{*} m^{|\omega|-1}\Tilde{\gamma}_\omega=\frac{1}{m}\sum_{\emptyset \neq \omega\subseteq 1{:}s }  C^{|\omega|}_{*} m^{|\omega|}\prod_{j\in\omega}\Tilde{\Gamma}_j\leq \frac{1}{m}\prod_{j=1}^s \Big(1+C_*\Tilde{\Gamma}_j m\Big). 
        \end{align*}
 \end{proof}
\begin{remark}\label{rmk:tractability}
    In settings where $s$ increases unboundedly while $\Vert f\Vert_q$ stays bounded, \cite[Lemma 3]{HICKERNELL2003286} shows if
    \begin{equation}\label{eqn:tractable}
      \lim_{s\to\infty}\sum_{j=1}^s\Tilde{\Gamma}_j<\infty,  
    \end{equation}
    then for any $\xi>0$, we can find $C_\xi$ independent of $s$ so that
    $$\prod_{j=1}^s \Big(1+C_*\Tilde{\Gamma}_j m\Big)\leq C_\xi b^{\xi m}$$
    and 
    $$\e \Big|\frac{1}{n}\sum_{i=0}^{n-1}f^w(\bsu_i)-\int_{\real^s} f(\bsx) \prod_{j=1}^s \varphi(x_j)\rd \bsx\Big|^2\leq C_\xi \Vert f\Vert^2_q b^{-(1+2\alpha-\xi)m}.$$

    For instance, when $\Gamma_j=O(j^{-\rho})$ for $\rho>2/q$ and $t_j=O(\log_b(j))$ as in the case of the Sobol' sequence and the Niederreiter sequence, 
    $$\Tilde{\Gamma}_j=O( j^{-2\rho+1+2\alpha} \theta^{-(1+\epsilon)(2\alpha-1+2/q) }_j+j^{1+2\alpha}\theta^{1-2\alpha}_j).$$
    By setting $\theta_j=\theta_0j^{-\rho q}$ with $\theta_0\in (0,1/2]$, a straightforward calculation shows
    $$\Tilde{\Gamma}_j=O( j^{-1-2(\rho q+1)(\alpha^*-\alpha)+\epsilon\rho q(2\alpha-1+2/q) }) \text{ for } \alpha^*=\frac{\rho q-2}{2(\rho q+1)}.$$
   It follows that the mean squared error of $\hat{\mu}$ converges at a dimension-independent rate arbitrarily close to $O(n^{-1-2\alpha^*})$ with the above choice of $\theta_j$.
\end{remark}


\section{Numerical experiments}\label{sec:numer}
To test our method, we consider standard Gaussian integrals with $\varphi(x)=\exp(-x^2/2)/\sqrt{2\pi}$, which satisfies Assumption~\ref{assump:rho}. In the experiments, we fix $\eta$ to be $\teta_1(u)=2^{-3} u^{-2}\exp(2-u^{-1})$. To compare the performance of $\hat{\mu}$ sharing the form~\eqref{eqn:muhat} under different choices of $T_j$, we fix $\bsu_0,\dots,\bsu_{n-1}$ to be linearly scrambled base-$2$ digital nets with direction numbers from \cite{joe:kuo:2008}. Each root mean squared error (RMSE) of $\hat{\mu}$ is estimated from $30$ independent runs.

%\subsection{Performance on simple test functions}

We consider test functions of the form
$$f(\bsx)=\prod_{j=1}^s \left( 1+j^{-2} g(x_j) \right) \text{ for } g(x_j)=\frac{\exp(M x^2_j)}{\sqrt{1-2M}}-1.$$
We require $M<0.5$ so that $f\in L^1(\real^s,\varphi)$.
Because $\int_\real g(x)\varphi(x)\rd x=0$, a straightforward calculation using \eqref{eqn:anovadecomp} shows $f_{v}=\prod_{j\in v} j^{-2}g(x_j)$. Because $ g \in W^{1,q,\infty}_{\mix}(\real,\varphi)$ for $q\in (1,1/2M)$, $f\in W^{1,q,\infty}_{\mix}(\real^s,\varphi)$ and 
\begin{align*}
    \Vert f_v\Vert_{W^{1,q,\infty}_{\mix}(\real^s,\varphi)}=&\Big(\prod_{j\in v} j^{-2}\Big)\Big(\sum_{v'\subseteq v}  \Vert \partial^{v'} \prod_{j\in v}g(x_j)\Vert^q_{L^{q,\infty}(\real,\varphi)}\Big)^{1/q}\\
    =&\Big(\prod_{j\in v} j^{-2}\Big)\Vert g\Vert^{|v|}_{W^{1,q,\infty}_{\mix}(\real,\varphi)} .
\end{align*}
Since $f_\emptyset=1$, $\Vert f \Vert_q= \sup_{v\subseteq 1{:}s}\Vert f_v \Vert_{W^{1,q,\infty}_{\mix}(\real^s,\varphi)}\geq 1$ and 
\begin{equation*}
 \gamma_v=\Vert f\Vert^{-1}_{q} \Vert f_v \Vert_{W^{1,q,\infty}_{\mix}(\real^s,\varphi)}\leq \prod_{j\in v} \Gamma_j \text{ for }  \Gamma_j=j^{-2}\Vert g\Vert_{W^{1,q,\infty}_{\mix}(\real,\varphi)}.   
\end{equation*}
We can therefore infer the convergence of our method from Corollary~\ref{cor:qmcvar}. 


Figure~\ref{fig:m00_comparison} shows the simulation results for $M=0$ and $s=5$ or $30$. In this case, $f(\bsx)=1$ and the usual inversion method with $T_j(u)=\Phi^{-1}(u)$ integrates $f$ exactly. We compare the performance of four options for $T_j$: 
\begin{itemize}
    \item Option 1: $T_j(u)$ in \eqref{eqn:Tj} with $\theta_j=0.1$.
    \item Option 2: $T_j(u)$ in \eqref{eqn:Tj} with $\theta_j=0.1/j^2$.
    \item Option 3: $T_j(u)=-\cot(\pi u)$.
    \item Option 4: $T_j(u)=a u-a(1-u)$.
\end{itemize}
Option 3 is the inverse CDF of a Cauchy distribution (also called the M{\"o}bius-transfor-mation in \cite{suzuki2025mobius}). Option 4 is the truncation method with $a=\sqrt{2\log n}$ suggested by \cite[Theorem 1b]{nuyens2023scaled} (we also tested $a=2\sqrt{\log n}$ suggested by \cite{dick2018optimal} and observed larger RMSEs). When $s=5$, all methods except Option 4 achieve a nearly $O(n^{-1})$ convergence rate, with Option 2 performing slightly better than Option 1 and 3. Option 4 seems already suffering from the dimensionality. When $s=30$, all methods are suffering from the high dimensionality, with Option 2 still maintaining a convergence rate close to $O(n^{-0.75})$. 

\begin{figure}[h!] % The 'h!' option tells LaTeX to "place the figure HERE if possible"
    \centering % Centers the entire figure content on the page
    
    % First subfigure (M=0.0)
    \begin{subfigure}[b]{0.48\textwidth} % 'b' aligns the subfigures by their bottom edge
        \centering
        \includegraphics[width=\textwidth]{figure/RMSE_s5_M0.0.png}
        \caption{$s=5, M=0$}
        \label{fig:s5_m00} % A label for referencing this subfigure
    \end{subfigure}
    \hfill % Adds a horizontal flexible space between the two subfigures
    % Second subfigure (M=0.3)
    \begin{subfigure}[b]{0.48\textwidth} % Set width to < 0.5\textwidth to fit on one line
        \centering
        \includegraphics[width=\textwidth]{figure/RMSE_s30_M0.0.png}
        \caption{$s=30, M=0$}
        \label{fig:s30_m00}
    \end{subfigure}
    
    \caption{Comparison of RMSEs for $M=0.0$ with $s=5$ or $30$. The first four legend labels correspond to Options 1-4, and the two reference lines are  proportional to $n^{-1/2}$ and $n^{-1}$, respectively. }
    \label{fig:m00_comparison} % A label for referencing the entire figure
\end{figure}

Figure~\ref{fig:m03_comparison} shows the simulation results for $M=0.3$ and $s=5$ or $30$. In this case, $f\notin L^2(\real^s,\varphi)$ and the plain Monte Carlo method for estimating $\mu$ has an infinite variance. We again compare the performance of Options 1-4. We follow \cite[Theorem 1b]{nuyens2023scaled} and set $a=\sqrt{5\log n}$ in Option 4. In addition to the previous four options, we also compare the usual inversion method (without importance sampling):
\begin{itemize}
    \item Option 5: $T_j(u)=\Phi^{-1}(u).$
\end{itemize}
By Remark~\ref{rem:rqmcrate}, the asymptotic convergence rate of Option 5 is close to $O(n^{-0.4})$, consistent with the simulation results. Options 1-3 perform similarly to the $M=0$ case, indicating that all of them are capable of handling the severe boundary growth.

\begin{figure}[h!] % The 'h!' option tells LaTeX to "place the figure HERE if possible"
    \centering % Centers the entire figure content on the page
    
    % First subfigure (M=0.0)
    \begin{subfigure}[b]{0.48\textwidth} % 'b' aligns the subfigures by their bottom edge
        \centering
        \includegraphics[width=\textwidth]{figure/RMSE_s5_M0.3.png}
        \caption{$s=5$, $M=0.3$}
        \label{fig:s5_m03} % A label for referencing this subfigure
    \end{subfigure}
    \hfill % Adds a horizontal flexible space between the two subfigures
    % Second subfigure (M=0.3)
    \begin{subfigure}[b]{0.48\textwidth} % Set width to < 0.5\textwidth to fit on one line
        \centering
        \includegraphics[width=\textwidth]{figure/RMSE_s30_M0.3.png}
        \caption{$s=30$, $M=0.3$}
        \label{fig:s30_m03}
    \end{subfigure}
    
    \caption{Comparison of RMSEs for $M=0.3$ with $s=5$ or $30$. The first five legend labels correspond to Options 1-5, and the two reference lines are  proportional to $n^{-1/2}$ and $n^{-1}$, respectively.  }
    \label{fig:m03_comparison} % A label for referencing the entire figure
\end{figure}

Our final experiment studies how the choice of $\theta_j$ affects the performance of boundary-damping IS. We set $M=0.25$ so that $f\in W^{1,q,\infty}_{\mix}(\real^s,\varphi)$ for $q\in (1,2)$.  Our analysis in Remark~\ref{rmk:tractability} suggests $\theta_j=\theta_0 j^{-4}$  with $\theta_0\in (0,1/2]$ should produce a near-optimal decay in $\Tilde{\Gamma}_j$. We therefore compare the following three options for $T_j$:
\begin{itemize}
    \item Option 6: $T_j(u)$ in \eqref{eqn:Tj} with $\theta_j=0.1/j^2$.
    \item Option 7: $T_j(u)$ in \eqref{eqn:Tj} with $\theta_j=0.1/j^4$.
    \item Option 8: $T_j(u)$ in \eqref{eqn:Tj} with $\theta_j=0.1/j^6$.
\end{itemize}
We also use the usual inversion method Option 5 as a baseline. The results for $s=128$ are shown in Figure~\ref{fig:m025}. We see Options 6-8 significantly outperform the baseline, indicating our boundary-damping IS successfully accelerates the convergence in high-dimensional settings. We also observe Option 7 performs the best among Options 6-8, confirming our prediction.

\begin{figure}[h!] % The 'h!' option tells LaTeX to "place the figure HERE if possible"
    \centering % Centers the entire figure content on the page
    \includegraphics[width=0.48\textwidth]{figure/RMSE_s128_M0.25.png}
    \caption{RMSEs for $M=0.25$ with $s=128$. The first four legend labels correspond to Options 5-8, and the two reference lines are  proportional to $n^{-1/2}$ and $n^{-1}$, respectively.  }
    \label{fig:m025} % A label for referencing the entire figure
\end{figure}


\section{Concluding remarks}\label{sec:conclusion}

In this paper, we have proposed a new class of importance sampling methods suitable for RQMC integration of functions with severe boundary growth. Both our theoretical bounds and simulation results demonstrate a significant improvement in the convergence rates compared to previous methods.

As a limitation, our analysis does not extend to the usual inversion method by taking the limit $\theta_j\to 0$ for $j\in 1{:}s$. One reason is that as $\theta\to 0$, $w_{\theta}$ converges to $1$ 
pointwise over $\itr$ but not in the Sobolev norm $W^{1,2}_{\mix}(\itr)$. It is interesting to ask whether we can establish the convergence rates without bounding the Sobolev norm and hence bridge our method with the inversion method. We leave this question for future research.

Another limitation is that the convergence rates established in Theorem~\ref{thm:qmcvar} do not improve when $f\in W^{1,q}_{\mix}(\real^s,\varphi)$ or $W^{1,q,\infty}_{\mix}(\real^s,\varphi)$ with $q>2$. It is worth studying how our method performs on integrands with mild or even no boundary growth. In particular, it is an open question whether our method can reproduce the dimension-independent convergence rates in \cite{nichols:2014} under the same assumptions.

In \cite{suzuki2025mobius}, the authors prove the M{\"o}bius-transformed trapezoidal rule achieve the optimal convergence rate in the one-dimensional $\rho$-weighted Sobolev spaces $W_{\rho}^{\alpha,2}(\real)$. We conjecture that trapezoidal rules combined with our boundary-damping IS can achieve the same convergence rates. A detailed analysis is beyond the scope of this paper and left for future study.


\section*{Appendix} 

This appendix contains the proofs of Lemmas~\ref{lem:embedding}, \ref{lem:anova}, and \ref{eqn:etabound}. We will need the following lemma.

\begin{lemma}\label{lem:tailbound}

    For $\varphi$ satisfying Assumption~\ref{assump:rho} and $\varepsilon\in (0,1)$,
    $$\int_\real \varphi(x)^\varepsilon\rd x\leq 2  c_\varepsilon^{-1}  \varepsilon^{-2}.$$
\end{lemma}
\begin{proof} Because $\varphi(x)=\varphi(-x)$, we know that $\Phi(0)=1/2$ and
    \begin{align*}  
        \int_\real \varphi(x)^\varepsilon\rd x = 2\int_{-\infty}^0 \frac{\varphi(x) }{\varphi(x)^{1-\varepsilon}} \rd x \leq 2 c_\varepsilon^{-1} \int_{-\infty}^0 \frac{1}{\Phi(x)^{(1-\varepsilon)(1+\varepsilon)}} \rd \Phi(x)= 2  c_\varepsilon^{-1}  \varepsilon^{-2} 2^{-\varepsilon^2}.
    \end{align*}
\end{proof}



\begin{proof}[Proof of Lemma \ref{lem:embedding}]
   Note that for any $\bsx \in \real^s$, the following inequality holds:
$$|f(\bsx)|^q \prod_{j=1}^s \varphi(x_j) \leq \varphi_\infty^s\Vert f \Vert_{L^{q,\infty}(\real^s,\varphi)}^q,$$
or equivalently
\begin{equation} \label{eq:pointwise_bound}
|f(\bsx)| \leq \varphi_\infty^{s/q}\Vert f \Vert_{L^{q,\infty}(\real^s,\varphi)} \left( \prod_{j=1}^s \varphi(x_j) \right)^{-1/q}.
\end{equation}
We thus have
\begin{align*}
\Vert f \Vert_{L^{q'}(\real^s,\varphi)}^{q'} &= \int_{\real^s} |f(\bsx)|^{q'} \prod_{j=1}^s \varphi(x_j) \rd\bsx \\
&\leq \varphi_\infty^{sq^{\prime}/q}\Vert f \Vert_{L^{q,\infty}(\real^s,\varphi)}^{q'}\int_{\real^s} \prod_{j=1}^s \varphi(x_j)^{1-q'/q} \rd\bsx\\
&=\varphi_\infty^{sq^{\prime}/q}\Vert f \Vert_{L^{q,\infty}(\real^s,\varphi)}^{q'}\left( \int_{\real} \varphi(x)^{1-q'/q} \rd x \right)^s.
\end{align*}
Denote $I_{\alpha} := \int_{\real} \varphi(x)^{1-\alpha} \rd x$ for $\alpha\in(0,1)$. Since $q>q'$, the exponent $1-q'/q\in(0,1)$. Lemma~\ref{lem:tailbound} shows that $I_{q'/q}<\infty$. Thus, we have
$$\Vert f \Vert_{L^{q'}(\real^s,\varphi)}^{q'} \leq (\varphi_\infty^{q^{\prime}/q} I_{q'/q})^s\Vert f \Vert_{L^{q,\infty}(\real^s,\varphi)}^{q'}.$$
Taking the $(1/q')$-th root of both sides gives the first claim
$$\Vert f \Vert_{L^{q'}(\real^s,\varphi)} \leq (\varphi_\infty^{q^{\prime}/q} I_{q'/q})^{s/q'} \Vert f \Vert_{L^{q,\infty}(\real^s,\varphi)}.$$
This proves the inequality with a constant $C_{q,q'} =  (\varphi_\infty^{q^{\prime}/q} I_{q'/q})^{1/q'}$.


For the second inequality, 
we first apply the result from the first part to each term $\partial^v f$,
$$\Vert \partial^v f \Vert_{L^{q'}(\real^s,\varphi)} \leq (\varphi_\infty^{q^{\prime}/q} I_{q'/q})^{s/q'} \Vert \partial^v f \Vert_{L^{q,\infty}(\real^s,\varphi)}.$$
Substituting this into the Sobolev norm definition gives
\begin{align*}
\Vert f \Vert_{W^{1,q'}_{\mix}(\real^s,\varphi)} &= \left(\sum_{v\subseteq \{1,\dots,s\}} \Vert \partial^v f \Vert^{q'}_{L^{q'}(\real^s,\varphi)}\right)^{1/q'}\\&\leq (\varphi_\infty^{q^{\prime}/q} I_{q'/q})^{s/q'}\left(\sum_{v\subseteq \{1,\dots,s\}} \Vert \partial^v f \Vert^{q'}_{L^{q,\infty}(\real^s,\varphi)} \right)^{1/q'}.
\end{align*}
Now, let $\bsb$ be a vector in $\mathbb{R}^{2^s}$ with components $b_v = \Vert \partial^v f \Vert_{L^{q,\infty}(\real^s,\varphi)}$. The expression above is $\varphi_\infty(I_{q'/q})^{s/q'} \Vert \bsb \Vert_{q'}$. For a finite-dimensional vector space, we know that for $q > q'$, the $\ell_{q'}$ norm is bounded by the $\ell_q$ norm. The dimension of our vector space is the number of subsets $v$, which is $d = 2^s$. The inequality is
$$\Vert \bsb \Vert_{q'} \leq d^{(1/q' - 1/q)} \Vert \bsb \Vert_q = 2^{s(q-q')/(qq')} \Vert \bsb \Vert_q.$$
Applying this inequality gives
$$\Vert f \Vert_{W^{1,q'}_{\mix}(\real^s,\varphi)} \leq (\varphi_\infty^{q^{\prime}/q} I_{q'/q})^{s/q'} \cdot 2^{s(q-q')/(qq')} \Vert f \Vert_{W^{1,q,\infty}_{\mix}(\real^s,\varphi)}.$$
Taking $\tilde C_{q,q'} = (\varphi_\infty^{q^{\prime}/q} I_{q'/q})^{1/q'} 2^{(q-q')/(qq')}$ completes the proof.
\end{proof}

\begin{proof}[Proof of Lemma \ref{lem:anova}]

For the $W^{1,q}_{\mix}(\real^s,\varphi)$ case, we first show that $P_j$ is a contraction on $W^{1,q}_{\mix}(\real^s,\varphi)$. Let $g \in W^{1,q}_{\mix}(\real^s,\varphi)$. Note that  $\partial^u P_j g = P_j \partial^u g$ if $j \notin u$, and $\partial^u P_j g = 0$ if $j \in u$,
\begin{align*}
\Vert P_j g \Vert^q_{W^{1,q}_{\mix}(\real^s,\varphi)} &= \sum_{u \subseteq 1{:}s} \Vert \partial^u (P_j g) \Vert^q_{L^q(\real^s,\varphi)} \\
&= \sum_{u \subseteq 1{:}s, j \notin u} \Vert P_j (\partial^u g) \Vert^q_{L^q(\real^s,\varphi)}.
\end{align*}
By Jensen's inequality, for any function $h$, we have $|P_j(h)(\bsx)|^q \le P_j(|h|^q)(\bsx)$. Integrating this over $\real^s$ with the weight $\prod_k \varphi(x_k)$ shows that $\Vert P_j h \Vert_{L^q(\real^s,\varphi)} \leq \Vert h \Vert_{L^q(\real^s,\varphi)}$. Thus, $P_j$ is a contraction on $L^q(\real^s,\varphi)$. Applying this, we get
\begin{align*}
\Vert P_j g \Vert^q_{W^{1,q}_{\mix}(\real^s,\varphi)} &\leq \sum_{u \subseteq 1{:}s, j \notin u} \Vert \partial^u g \Vert^q_{L^q(\real^s,\varphi)} \\
&\leq \sum_{u \subseteq 1{:}s} \Vert \partial^u g \Vert^q_{L^q(\real^s,\varphi)} = \Vert g \Vert^q_{W^{1,q}_{\mix}(\real^s,\varphi)}.
\end{align*}
This shows $\Vert P_j \Vert_{W^{1,q}_{\mix} \to W^{1,q}_{\mix}} \leq 1$. By the triangle inequality, the operator $(I-P_j)$ is also bounded: $\Vert I-P_j \Vert \leq \Vert I \Vert + \Vert P_j \Vert \leq 2$.

Since $f_v = (\prod_{j\in v} (I-P_j)) P_{1{:}s\setminus v} f$ is a composition of bounded linear operators applied to $f$, and $f \in W^{1,q}_{\mix}(\real^s,\varphi)$, it follows that $f_v \in W^{1,q}_{\mix}(\real^s,\varphi)$.

The $W^{1,q,\infty}_{\mix}(\real^s,\varphi)$ case requires Assumption~\ref{assump:rho}. We first show that $P_j$ is bounded on $L^{q,\infty}(\real^s,\varphi)$. For any $h \in L^{q,\infty}(\real^s,\varphi)$, we have the pointwise bound $|h(\bsx)| \leq \varphi_\infty^{s/q} \Vert h \Vert_{L^{q,\infty}} (\prod_k \varphi(x_k))^{-1/q}$. As a result,
\begin{align*}
|P_j h(\bsx)| &= \left| \int_{\real} h(\bsx) \varphi(x_j) \rd x_j \right| 
\leq \int_{\real} |h(\bsx)| \varphi(x_j) \rd x_j \\
&\leq \varphi_\infty^{s/q}\Vert h \Vert_{L^{q,\infty}(\real^s,\varphi)}\int_{\real}  \Big(\prod_{k=1}^s \varphi(x_k)\Big)^{-1/q}\varphi(x_j) \rd x_j \\
&= \varphi_\infty^{s/q}\Vert h \Vert_{L^{q,\infty}(\real^s,\varphi)}I_{1/q} \left(\prod_{k \neq j} \varphi(x_k)\right)^{-1/q}  ,
\end{align*}
where $I_{1/q}=\int_{\real}\varphi(y)^{1-1/q} \rd y<\infty$ by Lemma~\ref{lem:tailbound}.
Now we use this to bound the $L^{q,\infty}(\real^s,\varphi)$ norm of $P_j h$, yielding
\begin{align*}
\Vert P_j h \Vert^q_{L^{q,\infty}(\real^s,\varphi)} &= \sup_{\bsx \in \real^s} |P_j h(\bsx)|^q \prod_{k=1}^s \frac{\varphi(x_k)}{\varphi_\infty} \\
&\leq \sup_{\bsx \in \real^s}  \Vert h \Vert^q_{L^{q,\infty}(\real^s,\varphi)} I_{1/q}^q\left(\prod_{k \neq j} \varphi(x_k)\right)^{-1} \prod_{k=1}^s \varphi(x_k) \\
&= \varphi_\infty\Vert h \Vert^q_{L^{q,\infty}(\real^s,\varphi)} I_{1/q}^q.
\end{align*}
By Assumption~\ref{assump:rho}, both the integral and the supremum are finite. Let $C_j = \varphi_\infty^{1/q}I_{1/q}$. Since $\Vert P_j h \Vert_{L^{q,\infty}(\real^s,\varphi)} \leq C_j \Vert h \Vert_{L^{q,\infty}(\real^s,\varphi)}$, $P_j$ is bounded on $L^{q,\infty}(\real^s, \varphi)$.
Using the same reasoning as in the first part, for any $g \in W^{1,q,\infty}_{\mix}(\real^s,\varphi)$,
\begin{align*}
\Vert P_j g \Vert^q_{W^{1,q,\infty}_{\mix}(\real^s,\varphi)} &= \sum_{u \subseteq 1{:}s, j \notin u} \Vert P_j (\partial^u g) \Vert^q_{L^{q,\infty}(\real^s,\varphi)} \\
&\leq \sum_{u \subseteq 1{:}s, j \notin u} C_j^q \Vert \partial^u g \Vert^q_{L^{q,\infty}(\real^s,\varphi)} \leq C_j^q \Vert g \Vert^q_{W^{1,q,\infty}_{\mix}(\real^s,\varphi)}.
\end{align*}
Thus, $P_j$ is a bounded operator on $W^{1,q,\infty}_{\mix}(\real^s,\varphi)$. Consequently, $I-P_j$ is also bounded. Since $f_v$ is formed by applying these bounded operators to $f$, and $f \in W^{1,q,\infty}_{\mix}(\real^s,\varphi)$, we conclude that $f_v \in W^{1,q,\infty}_{\mix}(\real^s,\varphi)$.
\end{proof}

\begin{proof}[Proof of Lemma \ref{eqn:etabound}]
Note that for $u \in (0,1/2]$, $ \eta_p^{\prime}(u) = (pu^{-p} - p - 1)u^{-1} \eta_p(u) \le pu^{-p-1} \eta_p(u)$. %\Du{Is this always positive?} \Pan{It is when $p\in (0,1]$. We also need to prove it.} \Du{If $u = 1/4$ and $p\rightarrow 0$, it seems that $4^p p - p - 1$ will be negative.} \Pan{Good point. I think the condition should be $p\geq 1$.}
    Thus, we have $  \eta_p(u) \ge 0$ and $ \eta_p(u) \ge \frac{u^{p+1}}{p}\eta_p^{\prime}(u)$, which implies that $ \eta_p$ is increasing and, moreover,
    $$
    \int_0^u  \eta_p(t)\rd t = 2^{-p-2}\frac{1}{p}\exp(2^p - u^{-p}) = \frac{u^{p+1}}{p} \eta_p(u).
    $$
    This proves the desired results.
\end{proof}


%\section*{Acknowledgments}



\bibliographystyle{siamplain}
\bibliography{qmc}
\end{document}
